% !TEX root = ../main.tex
\chapter{구현과 실증 결과}
\label{ch:results}

\ref{ch:methodology}장에서는 점수를 어떻게 만드는지, 어떤 점검을 하는지, 불확실성(결과가 얼마나 흔들릴 수 있는지)을 어떻게 수로 나타내는지, 실행 시간을 어떻게 재는지 설명했어요. 이 장에서는 그 설계를 아비트럼 원(Arbitrum One)에서 하나의 사례 연구로 실제로 해 봐요. 블록체인 하나와 관찰 기간 하나를 써요. 같은 기간 안의 비교는 모두 매칭 코호트(모든 점수가 똑같이 점수를 매기는 지갑 묶음) 하나에서 해요. 시간 홀드아웃(나중 기간을 떼어 두고 맞혀 보는 시험)에는 스펜더(spender, 허락을 받아 대신 쓰는 쪽) 코호트와 트레이더(거래하는 사람) 코호트를 써요. 점(노드) 개수에 따른 민감도(조건을 바꿀 때 결과가 얼마나 움직이는지)를 볼 때는 확장 주소 풀(더 큰 주소 모음) 하나를 써요. 그리고 \ref{sec:holdout-results}절 끝에는 2026년 6월부터 8월까지의 자료로 한 등록 재현(계획을 먼저 공개해 두고 새 자료로 다시 해 보기)이 있어요. 모든 표는 \ref{sec:reproducibility}절의 평가 요약에서 만들었어요. 켄달의 타우($\tau$, 두 줄 세우기가 얼마나 같은 순서인지 나타내는 수)에는 언제나 신뢰구간(믿을 만한 범위)을 붙였어요. 결과마다 모든 점수를 함께 보고해요. 페이지랭크(PageRank) 점수는 저마다, 그 점수가 매끄럽게 다듬는 원시 차수(화살표를 그냥 센 수)와 견주어 읽어요. 엔도스랭크(EndorseRank)와 AWP(적응형 가중 페이지랭크) 사이의 짝지은 차이(같은 지갑들로 두 점수를 견준 차이)는 부록~\ref{ch:appendix-er-awp}에 있어요.

\dissertationSubheading[sec:dataset]{사례 연구 배경: 아비트럼 원}

아비트럼 원을 고른 까닭은 세 가지예요. 아비트럼 원은 많이 쓰이는 이더리움 롤업(거래를 묶어 이더리움에 올리는 레이어 2 블록체인)이고, 오래 다듬어진 디파이(DeFi, 탈중앙 금융) 서비스들이 있어요(\ref{sec:ethereum-l2}절). 또 아비트럼 원의 로그(계약이 남기는 기록)는 누구나 쓸 수 있는 빅쿼리(BigQuery) 공개 자료에 들어 있고, 이 연구의 처리 과정(파이프라인)이 바로 그 자료를 읽어요. 그리고 아비트럼 원의 증거금 거래(돈을 빌려 더 크게 하는 거래) 서비스 가운데 하나인 GMX~V2가 처음 고를 지갑들(씨앗 지갑 묶음)을 줘요. 관찰 기간은 2025년 12월 1일부터 2026년 5월 31일까지예요. 이 기간에는 시장 상황이 여러 번 바뀌었어요. 그리고 이 기간 안에서는 어떤 이벤트의 시간 무게(오래된 일일수록 가볍게 세는 무게)가 최근 이벤트 무게의 약 절반 아래로 떨어지는 일이 없어요(그림~\ref{fig:logistic-decay-methods}).

\dissertationSubsubheading[sub:cohorts]{코호트}

주소 묶음 세 가지를 써요.

\begin{itemize}
\item 매칭 코호트($n=5{,}521$). 관찰 기간에 포지션(거래하려고 열어 둔 자리)을 적어도 세 번 닫은 GMX~V2 지갑들이에요. 이 묶음은 허락(allowance, 토큰을 얼마까지 대신 써도 된다는 허락)이나 송금(transfer) 자료를 하나라도 꺼내기 전에 GMX 자료로 정해 두었어요(\ref{sec:data-sources}절). 그래서 어떤 지갑이 이 묶음에 드는지는 두 그래프(점과 화살표로 이은 그림) 어느 쪽에도 달려 있지 않아요. 코호트 지갑 가운데 $966$개에는 허락 화살표(엣지)가 하나도 닿지 않고, $381$개에는 송금 화살표가 하나도 닿지 않아요. $370$개에는 어느 쪽 화살표도 닿지 않아요. 어떤 그래프에 없는 지갑은 그 그래프를 쓰는 방법에서 0점을 받아요. 같은 기간 일치도(점수가 같은 기간의 지표와 얼마나 같은 순서인지), 차이 검사, 벤치마크(속도 재기)는 모두 이 코호트에서 계산했어요.
\item 홀드아웃 스펜더 코호트($n=1{,}335$). 꺼낸 허락 자료에서 $t_1$(2026년 2월 28일 23:59:59~UTC)에 0보다 큰 마지막 허락(가장 최근 허락)을 적어도 하나 가진 스펜더는 모두 여기에 들어요. 대부분은 계약(스마트 계약, 블록체인 위에서 저절로 실행되는 프로그램)이에요. 2026년 9월 말에 아비트럼 노드(블록체인 기록을 가진 컴퓨터)에서 읽어 보니(블록 510{,}408{,}687), 주소 1,335개 가운데 1,174개에는 계약 코드가 있었고 161개는 외부 소유 계정(EOA, 개인 키로 움직이는 계정)이었어요. 그 가운데 하나는 EIP-7702에 따라 계약 코드에 일을 맡겨요. 엔도스랭크가 가장 높게 매긴 스펜더 100개는 모두 계약이에요. 3월부터 5월까지 새 허락 쌍(새로 생긴 토큰 주인--스펜더 허락)을 얻은 스펜더 222개 가운데 221개도 계약이에요. 점수를 매긴 기간을 빅쿼리로 점검한 결과도, 판단할 수 있는 곳에서는 같은 답을 내요. 이벤트 로그를 남긴 스펜더는 모두 코드를 가지고 있기 때문이에요. 그 기간에 로그도 남기지 않고 거래도 보내지 않은 스펜더 686개 가운데 530개는 계약이에요. 그 가운데에는 자기 이벤트를 다른 계약이 대신 남기는 라우터(거래를 이어 주는 계약)들이 있어요. 파이프라인 설정에 이름이 적힌 GMX 거래소 라우터가 그런 예예요. 나머지 156개는 외부 소유 계정이에요. 이 지갑들의 점수는 2026년 3월 1일부터 5월 31일까지 처음 나타난 주인--스펜더 허락 쌍과 송금자(송금을 보낸 주소)에 견주어 시험해요(\ref{sec:holdout-protocol}절).
\item 확장 지갑 풀($N=99{,}080$). 매칭 코호트에, 이 블록체인의 달마다 ERC-20(토큰 공통 규칙) 활동에서 뽑은 주소 $93{,}559$개를 더한 거예요. 이 풀은 실행 시간의 점 개수 민감도와, 표본을 어떻게 정하느냐에 따른 민감도를 볼 때만 써요(\ref{sub:expanded-pool}절). 더한 주소들은 화살표를 하나도 가져오지 않아요. 그래서 풀 전체에서 풀이 프로그램(solver)이 푸는 그래프는 코호트 그래프와 같아요.
\end{itemize}

그래프 크기는 표~\ref{tab:benchmark-runtime}에 있어요. 같은 지갑들에서 보증 그래프(믿고 맡김을 나타낸 그래프)의 화살표 수는 송금 그래프의 약 9분의 1이에요. 엔도스랭크는 토큰, 주인(토큰 주인), 스펜더의 한 조합마다 마지막 허락 하나만 남기지만, AWP는 송금이 하나 있을 때마다 항을 하나씩 더하기 때문이에요. 표본 뽑기는 둘 다 같아요. 그래서 점수 사이의 차이는 화살표와 재시작(다시 출발하는 방식)에서 비롯돼요.

\dissertationSubsubheading[sub:implementation]{구현}

모든 알고리즘은 파이썬으로 짠 페이지랭크 구현 하나 위에서 돌아가요(감쇠 계수 $d=0.85$, 허용 오차 $\varepsilon=10^{-8}$, 최대 반복 횟수 $I_{\max}=300$). 엔도스랭크와 AWP는 화살표 목록과 재시작 벡터(어디서 다시 출발하는지 정한 값들)가 달라요. C-PR(결합 페이지랭크, 두 그래프를 함께 걷는 점수)은 두 화살표 목록을 점마다 섞어요(식~\eqref{eq:coupled-transition}). 한 번 돌릴 때마다 점수, 대리 지표(대신 재는 값)별 통계, 부트스트랩(다시 뽑기) 구간, 걸린 시간을 파일 하나에 적어요. 이 장의 표들은 2026년 10월 2일에 돌린 결과에서 나왔어요. 다만 등록 재현의 표들은 10월 1일에 평가한 결과예요. 돌린 목록, 저장소 판(리비전), 설정 파일, 요약 파일은 부록~\ref{ch:appendix-eval}에 있어요.

\dissertationSubheading[sec:benchmark-results]{계산 벤치마크(주장~C1)}

주장~C1은 같은 입력에서 각 점수를 계산하는 데 얼마나 드는지에 관한 거예요. \ref{sec:computational-benchmark}절의 시간 재는 방법을 따라, 표~\ref{tab:benchmark-runtime}은 매칭 코호트에서 알고리즘마다 다음 값들을 보고해요. 페이지랭크를 푸는 시간과 시간을 잰 다섯 번 되풀이에서의 표준편차(SD, 값이 흩어진 정도), 몸풀기로 한 번 푼 계산에서 가장 많이 쓴 메모리, 수렴(값이 더는 거의 바뀌지 않게 됨)까지의 반복(한 바퀴) 횟수, 평가에 쓴 부분 그래프의 크기예요.

% Real BigQuery data -- regenerate with:
%   2-Dissertation-Draft-Works/wallet-reputation-experiments/scripts/run_dissertation_eval.py --real --export-latex
% Generated by export_latex_results.py -- do not edit by hand unless re-export fails
\begin{table}[htbp]
\centering
\caption[Computational benchmark on the matched wallet cohort.]{Computational benchmark on the matched wallet cohort ($n=5{,}521$; AWP/EndorseRank runtime ratio 9.2$\times$). Runtime = mean wall-clock seconds over 5 timed PageRank solves after one untimed warm-up; SD over the same 5 runs. $|V|$, $|E|$ = nodes and edges of each method's evaluation subgraph. C-PR walks on the union of the allowance layer and the transfer layer, so $|E|$ is the sum of the two edge sets; S-PR runs on the transfer layer and its runtime includes the solve of the allowance layer that supplies its restart vector.}
\label{tab:benchmark-runtime}
\fitwidth{%
\begin{tabular}{lrrrrrr}
\toprule
Method & Runtime (s) & SD (s) & Peak mem.\ (MB) & Iters & $|V|$ & $|E|$ \\
\midrule
EndorseRank & 0.547 & 0.014 & 4.6 & 37.0 & 6{,}442 & 14{,}727 \\
AWP & 5.029 & 0.062 & 38.2 & 91.0 & 30{,}791 & 137{,}087 \\
\midrule
C-PR ($\lambda=0.5$) & 5.667 & 0.091 & 41.3 & 76.0 & 31{,}242 & 151{,}814 \\
S-PR & 5.600 & 0.020 & 39.1 & 145.0 & 30{,}791 & 137{,}087 \\
\bottomrule
\end{tabular}
}
\end{table}


시간을 잰 호출에서 대부분의 시간은 첫 반복 전에 쓰여요. 이때 파이썬 반복문이 화살표를 하나씩 돌면서, 화살표 목록으로 듬성듬성한 행렬(희소 행렬, 0이 아닌 칸이 적은 행렬)을 조립해요. 같은 화살표 목록으로 같은 호출을 따로 한 번 더 돌려 어디에 시간이 드는지 살펴보니(프로파일링), 실제 걸린 시간의 약 94\%가 행렬 조립에 들었어요(엔도스랭크는 $0.55$\,s 가운데 $0.52$\,s, AWP는 $5.17$\,s 가운데 $4.84$\,s). 거듭제곱 반복(같은 계산을 되풀이하기)은 약 1\%를 차지해요. 엔도스랭크의 $37$번 반복에 약 $6$\,ms, AWP의 $91$번 반복에 $67$\,ms가 걸려요. 나머지 시간은 대부분 점에 번호를 매기는 데 들고, AWP의 재시작 벡터를 만드는 데 약 $11$\,ms가 들어요. 조립과 반복 한 바퀴는 둘 다 $|E|$(화살표 수)에 비례해요(부록~\ref{sec:pr-sparsity}). 그래서 잰 시간의 비율 $9.2\times$는 화살표 수의 비율 $9.3\times$를 따라가요. 엔도스랭크는 수렴하는 데 필요한 반복 횟수도 더 적어요. 허락 그래프에는 자기 자신으로 돌아오는 화살표 하나 말고는 빙 돌아오는 길(순환)이 없어요. 그리고 화살표의 $99.0\%$는 나가는 화살표가 없는 스펜더에서 끝나요. 그래서 모든 길은 한 걸음 만에 멈추고, 걷는 확률(질량)은 재시작 분포(어디서 다시 출발하는지의 비율)를 거쳐 되돌아가요. 반면 송금 그래프에는 서로 오갈 수 있게 단단히 이어진 주소 $9{,}278$개의 중심부가 있고, 화살표의 $28.5\%$에는 반대 방향 화살표도 있어요. 그래서 걷기가 더 오래 돌고 돌아요. 이 차이는 진짜지만, 실제 걸린 시간에는 거의 보태지 않아요. 가장 많이 쓴 메모리도 마찬가지로 $|E|$에 따라 커져요. 돌릴 때마다 생기는 흔들림은 방법 사이의 차이보다 자릿수로 두 단계쯤(약 100배) 작아요(표준편차 $0.014$\,s와 $0.062$\,s).

혼합 점수 두 줄은 이 속도 결과가 어디까지 맞고 어디서 멈추는지 보여 줘요. C-PR은 두 화살표 집합을 합친 그래프(화살표 $151{,}814$개, 점 $31{,}242$개) 위를 걷고 $5.67$\,s가 걸려서, AWP보다 오래 걸려요. 반복 횟수는 더 적지만($76$ 대 $91$), 조립할 화살표가 더 많고, 호출 시간은 조립이 거의 다 차지해요. S-PR(보증 씨앗 페이지랭크)은 송금 층(그래프 한 겹) 위에서 돌고 $5.60$\,s가 걸려요. S-PR의 시간과 반복 횟수에는, 재시작 벡터를 대 주는 허락 층을 푸는 계산도 들어 있어요. $145$번의 반복은 그 계산의 $37$번에, 씨앗을 쓴 송금 걷기의 $108$번을 더한 거예요. 그러니까 싼 비용은 엔도스랭크만의 것이에요. 송금 층을 쓰는 알고리즘은, 송금 층만 쓰든 허락 정보와 함께 쓰든, 그 비용을 치러요.

이 결과에는 주의할 점이 두 가지 붙어요. 이 시간은 미리 만들어 둔 화살표 목록으로 풀이 프로그램을 부르는 시간만 쳐요. 처음부터 모두 새로 하면 읽어 들이기, 해독하기, 화살표 목록 만들기가 더해지는데, 이 일들은 모든 알고리즘이 전부 또는 일부 함께 써요. 그리고 더 빠른 방법이 꼭 더 좋은 방법은 아니에요. 그래서 실행 시간은 뒤에 나오는 일치도와 홀드아웃 결과에 비추어 읽어야 해요.

\dissertationSubsubheading[sub:runtime-scaling]{확장 풀에서 본 점 개수 민감도}

확장 풀에서 SHA256을 써서 언제나 같은 결과가 나오게 일부를 뽑으면(결정적 부분 표본 뽑기), 점 개수가 점점 커지는 세 단계가 생겨요. 표~\ref{tab:benchmark-scaling}는 각 단계의 실행 시간을, 그 단계에서 알고리즘마다 쓴 부분 그래프의 점과 화살표 수와 함께 보고해요. 코호트 안의 단계들($n\le 5{,}521$)은 부록~\ref{ch:appendix-eval}(표~\ref{tab:benchmark-scaling-incohort})에 있어요.

% Real BigQuery data -- regenerate with:
%   2-Dissertation-Draft-Works/wallet-reputation-experiments/scripts/run_dissertation_eval.py --real --export-latex
% Generated by export_latex_results.py -- do not edit by hand unless re-export fails
\begin{table}[htbp]
\centering
\caption[Runtime versus evaluation-set size on the expanded wallet pool.]{Runtime versus evaluation-set size on the expanded wallet pool ($N=99{,}080$; SHA256 subsampling seed \texttt{benchmark-tier2-v1}). Each stage keeps the edges that touch at least one sampled wallet, so $|V|$ and $|E|$ are reported per stage; the pool adds addresses but no additional extracted edges beyond the matched-cohort subgraph, hence the full-pool row reproduces the matched-cohort graph. Speedup = AWP/ER runtime ratio within the same row. Runtime = mean wall-clock seconds over 5 timed PageRank solves after one untimed warm-up; SD over the same 5 runs.}
\label{tab:benchmark-scaling}
\fitwidth{%
\begin{tabular}{rrrrrrrr}
\toprule
$n$ & ER (s) & AWP (s) & Speedup & ER $|V|$ & ER $|E|$ & AWP $|V|$ & AWP $|E|$ \\
\midrule
10{,}000 & 0.075 & 1.078 & 14.4$\times$ & 1{,}296 & 2{,}055 & 8{,}489 & 26{,}732 \\
50{,}000 & 0.342 & 3.414 & 10.0$\times$ & 4{,}440 & 9{,}190 & 19{,}990 & 92{,}887 \\
99{,}080 & 0.547 & 5.029 & 9.2$\times$ & 6{,}442 & 14{,}727 & 30{,}791 & 137{,}087 \\
\bottomrule
\end{tabular}
}
\end{table}


표~\ref{tab:benchmark-scaling}의 두 가지 특징은 \ref{sub:expanded-pool}절에서 나와요. 먼저, 풀 전체 줄은 표~\ref{tab:benchmark-runtime}의 매칭 코호트 줄, 그리고 표~\ref{tab:robustness-damping}의 $d=0.85$ 줄과 똑같아요. 일부러 그렇게 설계했어요. 더한 주소들의 로그는 하나도 읽지 않았으므로 $N=99{,}080$에서 만들어지는 부분 그래프는 코호트 부분 그래프이고, 서로 다른 그래프마다 한 세션에서 한 번씩만 시간을 재기 때문이에요. 다음으로, $|E|$는 규모 표 하나 안에서는 $n$과 함께 커지지만, 두 가지 표본 뽑기를 서로 견주면 그렇지 않아요. $n=10{,}000$인 풀 단계는 엔도스랭크 화살표($2{,}055$개)가 $n=1{,}000$인 코호트 안 단계($2{,}584$개, 표~\ref{tab:benchmark-scaling-incohort})보다 적어요. 풀의 약 6\%인 코호트 지갑들의 로그만 꺼냈기 때문이에요. 이 표가 받쳐 주는 것은 범위가 한정된 주장이에요. 화살표 수가 5배에서 7배까지 차이 나는 단계들에서, AWP는 엔도스랭크보다 $8.6$배에서 $14.4$배 오래 걸리고, 엔도스랭크는 $0.55$\,s 이하에 머물러요. 비율이 가장 큰 곳은 가장 작은 풀 단계예요. 여기서 AWP는 $97$번, 엔도스랭크는 $32$번 반복해요. 주소 수와 함께 화살표도 늘어날 때 두 알고리즘이 어떻게 변하는지는 보여 주지 못해요. 더 많은 화살표를 대 주기로 했던 두 번째 단계의 자료 꺼내기가 끝나지 않았기 때문이에요. 그림~\ref{fig:runtime-scaling}은 코호트 안 단계와 풀 단계를 로그 눈금 축 하나에 함께 그려요.

\begin{figure}[htbp]
\centering
\includegraphics[width=0.9\textwidth]{\figpath{runtime-scaling.pdf}}
\caption[{뽑은 지갑 수에 따른 페이지랭크의 실제 걸린 시간이에요($x$축은 로그 눈금이에요). 실선은 코호트 안 단계예요($n\le 5{,}521$, 표~\ref{tab:benchmark-scaling-incohort}). 점선은 확장 풀 단계예요($N=99{,}080$, 표~\ref{tab:benchmark-scaling}). 풀 단계는 주소를 더하지만 꺼낸 화살표는 하나도 더하지 않아요. 그래서 가장 오른쪽 풀 단계의 점은 매칭 코호트 전체에 있어요.}]{풀이 프로그램을 한 번 부르는 데 실제 걸린 시간을 뽑은 지갑 수에 견주어 그렸어요($x$축은 로그 눈금이에요). 실선은 코호트 안 단계예요(표~\ref{tab:benchmark-scaling-incohort}). 점선은 확장 풀 단계예요(표~\ref{tab:benchmark-scaling}). 풀 단계는 꺼낸 화살표를 하나도 더하지 않으므로, 가장 오른쪽 풀 단계의 점은 매칭 코호트 전체예요. 실행 시간은 $n$이 아니라 각 단계의 화살표 수를 따라가요.}
\label{fig:runtime-scaling}
\end{figure}

\dissertationSubheading[sec:same-window-alignment]{같은 기간 일치도(주장~C2)}

주장~C2는 각 점수가, 그 점수를 계산한 기간 안에서 무엇과 맞는지를 물어요. 표~\ref{tab:alignment-hybrid}은 송금 가족, 허락 가족, 시빌 안정성 가족(가짜 주소를 많이 만들어 속이는 시빌 공격에 덜 흔들리는 지표 묶음)이라는 지표 가족(비슷한 지표 묶음)마다, 모든 점수의 가족 평균(한 가족에 든 지표들의 평균) 켄달 $\tau$를 보고해요. 값마다 짝지은 부트스트랩으로 구한 95\% 신뢰구간을 붙였어요. 원시 차수는 이 표에 줄이 없어요. 이 기간에서는 원시 차수가 곧 대리 지표이기 때문이에요. 입차수(들어오는 화살표 수)는 송금 가족에, 받은 허락 수(허락 입차수)는 허락 가족에 들어요. 9월 14일에 정한 C-PR의 보조 규칙(\ref{sec:holdout-protocol}절)을 빼면, 아래의 비교는 어느 것도 미리 정한 것이 아니에요(\ref{ch:methodology}장).

% Real BigQuery data -- regenerate with:
%   2-Dissertation-Draft-Works/wallet-reputation-experiments/scripts/run_dissertation_eval.py --real --export-latex
% Generated by export_latex_results.py -- do not edit by hand unless re-export fails
\begin{table}[htbp]
\centering
\caption[Same-window family-mean Kendall $\tau$ with bootstrap intervals for every score of this thesis.]{Same-window family-mean Kendall $\tau$ with bootstrap intervals for every score of this thesis (matched cohort, $n=5{,}521$). The scores are EndorseRank and AWP, the coupled operator at five values of $\lambda$ and the endorsement-seeded walk; $\lambda=1$ is the allowance layer walked alone (latest allowances in raw amounts) and $\lambda=0$ the transfer layer walked alone (raw amounts discounted by the decay of AWP); both restart uniformly. The raw degrees have no row because they are proxies of the transfer and allowance families. CI = 95\% percentile bootstrap over wallets (400 paired resamples, seed 42).}
\label{tab:alignment-hybrid}
\fitwidth{%
\begin{tabular}{lccc}
\toprule
Method & Transfer & Allowance & Sybil \\
\midrule
EndorseRank & 0.330 [0.311, 0.350] & 0.375 [0.342, 0.401] & 0.231 [0.212, 0.250] \\
AWP & 0.612 [0.602, 0.622] & -0.024 [-0.043, -0.002] & 0.278 [0.264, 0.292] \\
\midrule
C-PR ($\lambda=1$) & 0.333 [0.314, 0.353] & 0.355 [0.322, 0.384] & 0.233 [0.214, 0.253] \\
C-PR ($\lambda=0$) & 0.534 [0.522, 0.547] & -0.030 [-0.045, -0.017] & 0.286 [0.274, 0.300] \\
C-PR ($\lambda=0.25$) & 0.531 [0.520, 0.545] & -0.009 [-0.020, 0.004] & 0.288 [0.276, 0.302] \\
C-PR ($\lambda=0.5$) & 0.527 [0.515, 0.540] & -0.001 [-0.013, 0.013] & 0.290 [0.279, 0.304] \\
C-PR ($\lambda=0.75$) & 0.519 [0.507, 0.533] & 0.006 [-0.007, 0.020] & 0.292 [0.281, 0.306] \\
S-PR & 0.494 [0.481, 0.508] & 0.116 [0.097, 0.134] & 0.220 [0.207, 0.235] \\
\bottomrule
\end{tabular}
}
\end{table}


점수들은 어떤 화살표 위를 걷느냐에 따라 나뉘어요. 송금 위를 걷는 점수들, 곧 AWP, $\lambda$가 $1$보다 작은 모든 C-PR, S-PR은 송금 가족에서 $0.494$에서 $0.612$까지 나오고, 허락 가족에서는 0 가까이에 머물러요. 허락 가족에서는 허락 걷기의 결과에서 다시 출발하는 S-PR만 $0.116$에 이르러요. 허락 위만 걷는 두 점수, 곧 엔도스랭크와 C-PR ($\lambda=1$)은 허락 가족에서 $0.375$와 $0.355$, 송금 가족에서 약 $0.33$이 나와요. 시빌 안정성 가족에서는 모든 점수가 $0.220$에서 $0.292$ 사이에 있어요. 이 일치의 일부는 처음부터 생기게 되어 있어요. 각 가족의 주요 대리 지표는 두 그래프 가운데 하나와 같은 이벤트에서 나와요. 입차수와 받은 양의 합은 송금에서, 받은 허락 수와 받은 허락 금액은 마지막 허락에서 나와요. 그리고 페이지랭크는 자기 그래프의 입차수를 매끄럽게 다듬은 값이에요. 그러니까 점수가 자기 화살표 종류의 대리 지표와 맞는 것은 의도한 구성개념 점검(점수가 재려는 생각을 정말 재는지 보는 점검)이고, 바깥 시험으로 칠 수는 없어요. \ref{sec:holdout-results}절의 홀드아웃은 같은 질문을 기간 밖에서 물어요.

부록~\ref{sec:proxy-family-detail}는 엔도스랭크와 AWP의 가족 평균을 대리 지표 하나하나로 나누어 보여 줘요. AWP의 계수(여기서는 켄달의 타우 값)가 가장 큰 것은 송금 입차수(송금을 보내온 서로 다른 주소 수)와의 계수예요($\tau=0.728$). 받은 양의 합과의 계수는 훨씬 낮아요($0.495$). AWP의 무게를 보면 예상할 수 있는 일이에요. $V$는 원래 단위(소수점 조정을 하지 않은 가장 작은 단위)로 10 이상인 송금을, 금액이 얼마든 거의 한 번으로 세요. 그래서 AWP는 들어온 송금의 금액 합보다 들어온 송금의 개수를 더 가깝게 따라가요. 그렇다고 AWP가 입차수와 똑같아지지도 않아요. 퍼뜨림(전파), 시간 무게, 활발히 보내는 주소에서 다시 출발하기가, 두 지갑씩 짝지어 본 순서 가운데 잴 수 있을 만큼의 몫을 바꿔요. 그리고 기간 밖에서 AWP는 새 송금자를 입차수보다 못하게 줄 세워요(\ref{sec:holdout-results}절). 시빌 안정성 가족에서 AWP는 활동 기간(처음 활동부터 마지막 활동까지)과 활동한 달 수를 따라가고, 들어온 상대 비율(서로 다른 보낸 주소 수를 받은 송금 수로 나눈 값)에서는 뚜렷하게 음수예요($-0.228$). AWP의 화살표 무게는 송금을 세요. 그래서 적은 수의 주소에게서 송금을 많이 받는 지갑이 높은 순위에 오르는데, 이 비율은 바로 그런 모습에 벌점을 줘요.

엔도스랭크의 계수는 동점이 정해요(\ref{sec:rank-divergence}절). 허락을 하나도 받지 않은 지갑들의 공통 값보다 위에 있는 코호트 지갑은 $131$개뿐이에요. 두 허락 대리 지표가 0보다 큰 지갑도 바로 이 $131$개이고, 엔도스랭크는 이 지갑들을 모두 공통 값보다 위로 올려요. 그래서 엔도스랭크의 두 허락 계수($0.376$과 $0.373$)는 대부분 같은 가르기를 재는 것이고, 동점이 있을 때 가능한 천장(얻을 수 있는 가장 큰 값)에 붙어 있어요($0.377$과 $0.378$). 동점이 없는 대리 지표와 견줄 때 천장은 약 $0.57$이에요. 그래서 가족이 다르면 엔도스랭크의 계수들은 같은 잣대 위에 있지 않아요. 송금 계수($0.330$)는 주로, 허락 그래프 안의 지갑이 그 밖의 지갑보다 송금도 더 많이 받는다는 것을 나타내요. 그리고 그래프 밖의 지갑들을 외톨이 점(화살표가 하나도 없는 점)으로 남겨 두면 두 계수가 모두 크게 바뀌어요(\ref{sec:rank-divergence}절). C-PR ($\lambda=1$)은 같은 허락 그래프를 금액 그대로 걷고, 역시 모든 점에서 고르게 다시 출발해요(균등 재시작). 그리고 모든 가족에서 엔도스랭크와의 차이가 $0.02$ 안에 머물러요.

연구 계획서에 있던 같은 기간 거래 가족들, 곧 청산(손실이 너무 커져 거래가 강제로 끝나는 일), 위험 반대(값이 클수록 손실이 적은 지표), GMX 거래 성공 가족(\ref{sub:gf-pr}절)은 표~\ref{tab:alignment-summary}에 보관용으로 남겨 두었어요. 모든 점수가 이 가족들과는 약하게만 맞아요. 가족 평균이 $0.24$를 넘는 것은 없어요. 그리고 위험 반대 지표 두 개의 방향을 바르게 잡고 나면(부록~\ref{sec:supplementary-summary}), 모든 점수가 위험 반대 가족에서 음수예요. 손실 합계와 가장 나쁜 닫기(손실이 가장 컸던 한 번)는 지갑이 거래를 많이 할수록 커져요. 손실 없이 닫은 비율은 그렇게 커지지 않는 비율인데, 여기서 점수들의 계수는 $-0.031$에서 $0.024$ 사이에 있어요.

주장~C2는 좁은 뜻에서 뒷받침돼요. 한 종류의 화살표 위를 걷는 점수는 모두 그 화살표 종류의 대리 지표와 맞고, 그 구간은 0에서 충분히 떨어져 있어요. 그리고 섞은 점수들은 두 층 가운데 더 빽빽한 층을 따라가요. 이 일치는 일부가 처음부터 생기게 되어 있고, 엔도스랭크의 경우에는 그래프 밖의 지갑에 0점을 주는 규칙이 이 일치를 정해요.

\dissertationSubsubheading[sec:coupled-results]{결합 연산자(RQ5)}

\ref{sub:coupled-graph}절에서 혼합 점수 두 가지를 정했어요. 결합 연산자 C-PR은 지갑마다 이동 행렬(전이 행렬)의 줄에서, 금액을 그대로 쓴 허락 층과 송금 층을 섞는 비율 $\lambda$로 섞어요. 보증 씨앗 걷기 S-PR은 송금 층을 그대로 쓰되, 허락 걷기의 결과에서 다시 출발해요. 표~\ref{tab:tau-diff-hybrid}는 이 기간에서 두 점수의 짝지은 대비(같은 표본에서 구한 두 값의 차이)를 보여 줘요. 먼저 \ref{sec:statistical-tests}절의 대비 여덟 개는 가족마다 앞서는 층과 견준 것이고, 그다음 네 개는 엔도스랭크와 AWP와 견준 거예요. RQ5(연구 질문 5)의 답은 \ref{sec:holdout-results}절의 기간 밖 시험이 정해요. 여기서는 한 층만 쓰는 점수들이 통과하도록 만들어진 점검들 위에서, 혼합 점수들이 어디쯤 있는지만 살펴봐요.

% Real BigQuery data -- regenerate with:
%   2-Dissertation-Draft-Works/wallet-reputation-experiments/scripts/run_dissertation_eval.py --real --export-latex
% Generated by export_latex_results.py -- do not edit by hand unless re-export fails
\begin{table}[htbp]
\centering
\caption[Same-window paired contrasts for the coupled operator (C-PR, $\lambda=0.5$) and the endorsement-seeded walk (S-PR) against the layer that leads each family walked alone, the allowance layer (C-PR ($\lambda=1$)) or the transfer layer (C-PR ($\lambda=0$)), and of C-PR against EndorseRank and AWP.]{Same-window paired contrasts for the coupled operator (C-PR, $\lambda=0.5$) and the endorsement-seeded walk (S-PR) against the layer that leads each family walked alone, the allowance layer (C-PR ($\lambda=1$)) or the transfer layer (C-PR ($\lambda=0$)), and of C-PR against EndorseRank and AWP (matched cohort, $n=5{,}521$). Family-mean $\tau$; the same wallet resamples are used for $a$ and $b$. The contrasts with the endpoints are a stricter check beside the secondary criterion fixed on 14~September~2026 (Section~\ref{sec:holdout-protocol}), which compares C-PR's point estimate with the leading layer's interval in Table~\ref{tab:alignment-hybrid}. An interval below 0 means $a$ is significantly below $b$; one above 0 means it exceeds it. CI = 95\% percentile bootstrap over wallets (400 paired resamples, seed 42).}
\label{tab:tau-diff-hybrid}
\fitwidth{%
\begin{tabular}{lccccc}
\toprule
Contrast ($a$ minus $b$) & $\tau_a$ & $\tau_b$ & $\Delta\tau$ & 95\% CI of $\Delta\tau$ & $P(\Delta\tau>0)$ \\
\midrule
\multicolumn{6}{l}{\emph{Against the two layers walked alone}} \\
C-PR transfer minus C-PR ($\lambda=0$) transfer & 0.527 & 0.534 & -0.007 & [-0.009, -0.005] & 0.000 \\
C-PR allowance minus C-PR ($\lambda=1$) allowance & -0.001 & 0.355 & -0.356 & [-0.388, -0.323] & 0.000 \\
C-PR sybil-stability minus C-PR ($\lambda=0$) sybil-stability & 0.290 & 0.286 & +0.004 & [0.003, 0.006] & 1.000 \\
C-PR sybil-stability minus C-PR ($\lambda=1$) sybil-stability & 0.290 & 0.233 & +0.057 & [0.043, 0.069] & 1.000 \\
S-PR transfer minus C-PR ($\lambda=0$) transfer & 0.494 & 0.534 & -0.040 & [-0.049, -0.030] & 0.000 \\
S-PR allowance minus C-PR ($\lambda=1$) allowance & 0.116 & 0.355 & -0.239 & [-0.263, -0.215] & 0.000 \\
S-PR sybil-stability minus C-PR ($\lambda=0$) sybil-stability & 0.220 & 0.286 & -0.065 & [-0.072, -0.058] & 0.000 \\
S-PR sybil-stability minus C-PR ($\lambda=1$) sybil-stability & 0.220 & 0.233 & -0.013 & [-0.025, -0.003] & 0.007 \\
\midrule
\multicolumn{6}{l}{\emph{Against EndorseRank and AWP}} \\
C-PR transfer minus AWP transfer & 0.527 & 0.612 & -0.085 & [-0.095, -0.075] & 0.000 \\
C-PR allowance minus EndorseRank allowance & -0.001 & 0.375 & -0.375 & [-0.404, -0.344] & 0.000 \\
C-PR sybil-stability minus AWP sybil-stability & 0.290 & 0.278 & +0.012 & [0.005, 0.020] & 1.000 \\
C-PR sybil-stability minus EndorseRank sybil-stability & 0.290 & 0.231 & +0.059 & [0.045, 0.072] & 1.000 \\
\bottomrule
\end{tabular}
}
\end{table}


C-PR은 송금 층에 허락 층의 작은 보정을 더한 것처럼 움직여요. 송금 가족 계수는 $\lambda=0.25$에서 $0.531$, $\lambda=0.75$에서 $0.519$로 떨어지는데, 언제나 C-PR ($\lambda=0$)의 $0.534$보다 낮고 그 값과 $0.015$ 안쪽으로 차이 나요. 주된 설정에서 그 걷기와의 짝지은 차이는 $-0.007$이에요(구간 $[-0.009,-0.005]$). 작지만 0에서 떨어져 있어요. AWP는 송금 가족을 둘 다보다 더 잘 줄 세워요. C-PR은 AWP보다 $0.085$ 낮아요($[-0.095,-0.075]$). 허락 가족에서 C-PR은 어떤 $\lambda$에서도 0과 구별할 수 없어요. 그리고 자기 허락 층보다 한참 낮고($\Delta\tau=-0.356$, 구간 $[-0.388,-0.323]$), 엔도스랭크보다도 낮아요($-0.375$, $[-0.404,-0.344]$). 까닭은 허락 층이 듬성듬성하다는 것(희소성)이에요. 송금 층은 화살표가 허락 층의 약 9배이고, 점은 $30{,}791$개 대 $6{,}442$개예요(표~\ref{tab:benchmark-runtime}). 그래서 합친 그래프의 점 대부분에는 나가는 허락 화살표가 없어요. 그런 점에서는 식~\eqref{eq:coupled-transition}의 섞는 규칙이 $\lambda$가 얼마든 송금 줄로 돌아가요. $\lambda$는 두 층 모두에 나가는 화살표가 있는, 수가 적은 점들에서만 힘을 써요. 걷기가 그 적은 점들을 벗어나면, 순간 이동(아무 점으로나 건너뛰어 다시 출발하기)할 때까지 송금 화살표를 따라 다녀요. $\lambda$를 $0.75$로 올려도 허락 계수는 100분의 1보다 적게 움직여요. 섞은 점수는 끝점인 $\lambda=1$에서만 허락 구성개념을 되찾아요. $\lambda=0.75$에서도 허락 계수는 아직 $0.006$이에요.

C-PR이 한 층만 걷는 두 점수를 모두 넘는 것은 시빌 보정 안정성 가족에서뿐이에요. \ref{sec:holdout-protocol}절의 보조 기준이 바로 이것을 요구했어요. C-PR은 $0.290$이고, C-PR ($\lambda=0$)은 $0.286$, C-PR ($\lambda=1$)은 $0.233$이에요. 두 짝지은 구간은 모두 0보다 위에 있어요($[0.003,0.006]$과 $[0.043,0.069]$). 이 가족에서 C-PR은 AWP($0.278$)와 엔도스랭크($0.231$)도 넘어요. 송금 걷기보다 나아진 정도는 1000분의 4예요. 그래서 통계로 보면 차이지만, 실제로 쓰는 데에서는 아무 차이도 아니에요.

S-PR은 모든 가족에서, 빌려 온 층보다 못해요. 송금 걷기를 허락 걷기에서 다시 출발시키면, 송금 계수는 $0.494$로 떨어지고(C-PR ($\lambda=0$)과 견주어 $\Delta\tau=-0.040$), 허락 계수는 $0.116$까지밖에 오르지 않아요(C-PR ($\lambda=1$)과 견주어 $\Delta\tau=-0.239$). 그리고 시빌 보정 안정성은 두 층보다 모두 낮아져요. 재시작 벡터는 순간 이동할 때마다 걷기가 어디에 내려앉는지를 바꿔요. 하지만 이런 밀도의 그래프에서는 걷기가 어디서 시간을 보내는지를 바꾸지 못해요. 이것은 \ref{sec:holdout-results}절이 기간 밖에서 얻은 결과와 짝을 이루는, 같은 기간의 결과예요.

9월 14일에 정한 대로, 보조 기준은 이렇게 요구해요. 송금 가족과 허락 가족에서는 C-PR의 점 추정값(하나의 값으로 낸 추정)이 더 나은 층의 95\% 구간 안에 있어야 해요. 그리고 시빌 보정 안정성에서는 C-PR이 두 층을 모두 넘어야 해요. $\lambda=0.5$에서 C-PR은 세 부분 가운데 두 부분을 통과하고, 세 번째에서 떨어져요. 송금 계수($0.527$)는 C-PR ($\lambda=0$)의 구간($[0.522,0.547]$) 안에 들고, 시빌 보정 안정성은 두 층을 모두 넘어요. 그런데 허락 계수($-0.001$)는 C-PR ($\lambda=1$)의 구간($[0.322,0.384]$)에서 한참 밖에 있어요. 그래서 보조 규칙은 통과하지 못해요. 표~\ref{tab:tau-diff-hybrid}의 짝지은 대비는 이 규칙 옆에 더 엄격한 점검으로 함께 보고하는데, 이 대비로 따지면 송금 부분도 통과하지 못했을 거예요. 짝지은 대비에서는 C-PR이 송금 층보다 조금이지만 구간이 0을 벗어날 만큼 낮기 때문이에요. 같은 기간 표들에서는 더 빽빽한 층이 섞은 점수를 좌우해요. 섞은 점수에 한 층만 쓴 점수들에는 없는 예측하는 힘이 조금이라도 있는지라는 남은 질문은, 홀드아웃이 다뤄요.

\dissertationSubheading[sec:rank-divergence]{순위 갈라짐(H1 / RQ1)}

가설~H1과 RQ1은 엔도스랭크와 AWP가 같은 지갑들을 서로 다른 순서로 세우는지를 다뤄요. 두 점수 벡터(점수를 늘어놓은 목록) 사이의 방법 간 켄달 $\tau_b$(동점을 따로 셈하는 켄달의 타우)는 $0.360$이고, 구간은 $[0.342,0.379]$예요(표~\ref{tab:tie-sensitivity}의 마지막 줄). 그러니까 이 구간에는 0도 1도 들어 있지 않아요. 두 순위는 서로 관련이 있지만, 서로 바꿔 쓸 수 있는 것과는 거리가 멀어요. 두 점수가 우연보다 더 잘 맞는 것은 주로, 엔도스랭크가 0점을 주는 지갑들이 AWP에서도 바닥 가까이에 있기 때문이에요(아래를 보세요). 지갑 쌍의 $67.3\%$가 엔도스랭크에서 동점이므로, 두 점수가 모두 순서를 매기는 쌍은 전체의 $32.7\%$뿐이에요. 그리고 그 쌍들 가운데 $18.6\%$는 순서가 어긋나는 쌍이에요.

점수 하나하나의 분포(주변 분포)는 그림~\ref{fig:score-distributions}에 있어요. 엔도스랭크의 분포는 동점 덩어리 하나가 거의 다 차지해요. 엔도스랭크는 코호트에서 서로 다른 값을 $99$개만 가져요. $4{,}424$개 지갑은 점수 $9.62\times10^{-5}$를 함께 가져요. 이 점수는 허락을 하나도 받지 않은 지갑이 순간 이동과, 매달린 점(나가는 화살표가 없는 점)에서 다시 나누어 준 몫으로 받는 점수예요. $966$개는 허락 화살표가 하나도 닿지 않아서 0점이에요. 공통 값보다 위에 있는 것은 $131$개뿐이에요. AWP에서는 동점이 드물어서, 지갑 쌍의 $0.5\%$예요. 점수 하나의 분포로는 두 점수가 같은 지갑들을 어떻게 줄 세우는지 보여 줄 수 없어요. 그래서 그림~\ref{fig:rank-scatter}은 두 순위를 서로 맞대어 그려요. 동점인 지갑은 $\tau_b$에서처럼 가운데 순위에 놓아요. 그래서 엔도스랭크의 동점 덩어리는 세로 기둥으로 나타나요. 엔도스랭크가 떼어 내는 $131$개 지갑은 AWP 순위의 대부분에 흩어져 있어요(AWP 순위의 중앙값 $3{,}096$, 사분위 범위(가운데 절반이 들어가는 범위) $1{,}520$--$4{,}850$). 그러니까 엔도스랭크의 선두는 AWP의 선두가 아니에요. 동점인 $4{,}424$개 지갑의 기둥은 AWP 범위의 거의 전체에 걸쳐 있어요(사분위 범위 $1{,}149$--$3{,}571$). 0점인 $966$개 지갑은 대부분 AWP 바닥 가까이에 있고(AWP 순위의 중앙값 $4{,}966$), 그 가운데 $370$개는 AWP에서도 0점이에요.

동점 덩어리는 두 규칙에서 생기는데, 그 가운데 하나는 이 파이프라인이 고른 것이에요. 공통 값에 있는 4,424개 지갑은 허락을 하나도 받지 않은 주인으로서 허락 그래프 안에 있어요. 그래서 페이지랭크는 이 지갑들에게 순간 이동 몫과 매달린 점의 몫만 주고, 다른 것은 주지 않아요. 0점인 966개 지갑은 아예 그래프에 없고, 파이프라인은 그래프에 없는 지갑에 0점을 줘요(\ref{sec:graph-construction}절). 외톨이 점으로 남겨 두었다면 이 지갑들도 공통 값을 받았을 거예요. 이 규칙이 엔도스랭크가 코호트를 줄 세우는 순서의 대부분을 정해요. 엔도스랭크가 동점으로 두지 않는 지갑 쌍 4.99백만 개(499만 개) 가운데 4.27백만 개(427만 개, 85.7\%)는 공통 값의 지갑과 0점 지갑을 맞세운 쌍이에요. 그래서 매칭 코호트에서 엔도스랭크는 주로 허락 그래프 안의 지갑과 그 밖의 지갑을 갈라놓아요. 관찰 기간에 지갑이 해 주었거나 받은 마지막 허락 가운데 0이 아닌 것이 적어도 하나 있으면, 그 지갑은 그래프 안에 있어요. 그래프 밖의 966개 가운데 802개는 그 기간에 허락을 한 번도 하지 않았어요. 나머지 164개는 허락을 했지만, 모든 토큰과 스펜더에서 마지막으로 한 허락이 허락 금액을 0으로 정했어요.

그래프만 보아서는 0점 규칙이 외톨이 점으로 두는 방법보다 나을 까닭이 없어요. 그래서 그래프에 없는 지갑들을 외톨이 점으로 남겨 두고 같은 기간 일치도를 다시 계산했어요. 이 점검은 위의 결과를 안 뒤에 했어요. 고르게 순간 이동하는 방식에서는 외톨이 점이 들어오는 화살표가 없는 다른 점과 똑같은 점수를 받아요. 그래서 966개 지갑은 4,424개 덩어리에 합쳐지고, 나머지 점수에는 모두 같은 수 하나가 곱해져요. AWP에서 그래프에 없는 381개 지갑도 같은 방법으로 다루었어요. AWP의 재시작 방식에서는 아무것도 보내지 않은 외톨이 지갑이 재시작 몫을 받지 못하므로, 어느 쪽이든 0점이에요. 결과는 표~\ref{tab:tie-sensitivity}에 있어요. AWP의 계수는 하나도 움직이지 않아요. 엔도스랭크의 허락 계수는 $0.990$으로 올라가요(구간 $[0.988,0.992]$). 이제 엔도스랭크와 허락 대리 지표가 코호트를 같은 곳에서 가르기 때문이에요. 송금 계수는 $-0.071$로 떨어지고($[-0.086,-0.056]$), 시빌 안정성 계수와 거래 성공 계수는 $-0.036$과 $-0.017$로 떨어져요. 두 점수 사이의 일치는 $0.360$에서 $-0.025$로 떨어져요($[-0.044,-0.003]$). 그러니까 송금, 안정성, 거래에서 엔도스랭크가 보인 같은 기간의 양수 계수와 AWP와의 일치는, 어떤 트레이더가 허락 그래프 안에 있는지를 재는 거예요.

% Post hoc checks -- regenerate with:
%   2-Dissertation-Draft-Works/wallet-reputation-experiments/scripts/run_supplementary_checks.py
% Generated by export_latex_results.py -- do not edit by hand unless re-export fails
\begin{table}[htbp]
\centering
\caption[Same-window family-mean Kendall $\tau$ with wallets missing from a graph scored zero or kept as isolated nodes.]{Same-window family-mean Kendall $\tau$ with wallets missing from a graph scored zero or kept as isolated nodes (matched cohort, $n=5{,}521$). The main analysis scores a wallet that is missing from a method's graph zero. The rule concerns 966 wallets for EndorseRank and 381 for AWP. Under EndorseRank's uniform restarts an isolated node gets the score of a node without in-edges. AWP restarts only at addresses that sent a transfer, so an isolated wallet keeps the score zero and AWP's coefficients are the same under both rules. The last row is the agreement between the two scores. The check was run after the results were known. CI = 95\% percentile bootstrap over wallets (400 paired resamples, seed 42).}
\label{tab:tie-sensitivity}
\fitwidth{%
\begin{tabular}{lcccc}
\toprule
 & \multicolumn{2}{c}{Missing wallets scored zero} & \multicolumn{2}{c}{Missing wallets as isolated nodes} \\
\cmidrule(lr){2-3} \cmidrule(lr){4-5}
Family & EndorseRank & AWP & EndorseRank & AWP \\
\midrule
Transfer & 0.330 [0.311, 0.350] & 0.612 [0.602, 0.622] & -0.071 [-0.086, -0.056] & 0.612 [0.602, 0.622] \\
Allowance & 0.375 [0.342, 0.401] & -0.024 [-0.043, -0.002] & 0.990 [0.988, 0.992] & -0.024 [-0.043, -0.002] \\
Sybil stability & 0.231 [0.212, 0.250] & 0.278 [0.264, 0.292] & -0.036 [-0.049, -0.022] & 0.278 [0.264, 0.292] \\
\addlinespace
Liquidation & 0.093 [0.078, 0.110] & 0.141 [0.129, 0.155] & 0.082 [0.069, 0.096] & 0.141 [0.129, 0.155] \\
Inverse risk & -0.055 [-0.071, -0.039] & -0.127 [-0.140, -0.115] & 0.006 [-0.015, 0.023] & -0.127 [-0.140, -0.115] \\
GMX trading success & 0.096 [0.080, 0.113] & 0.236 [0.223, 0.250] & -0.017 [-0.033, -0.003] & 0.236 [0.223, 0.250] \\
\midrule
EndorseRank vs.\ AWP & \multicolumn{2}{c}{0.360 [0.342, 0.379]} & \multicolumn{2}{c}{-0.025 [-0.044, -0.003]} \\
\bottomrule
\end{tabular}
}
\end{table}


\begin{figure}[htbp]
\centering
\includegraphics[width=0.9\textwidth]{\figpath{score-distributions.pdf}}
\caption[{매칭 코호트($n=5{,}521$)에서 엔도스랭크와 AWP 점수의 주변 분포예요. 점수가 0보다 큰 지갑을 $\log_{10}$ 눈금으로 그렸어요(엔도스랭크의 0점 지갑 $966$개와 AWP의 $381$개는 막대그래프에서 뺐어요).}]{매칭 코호트($n=5{,}521$)에서 엔도스랭크와 AWP 점수의 주변 분포예요. 점수가 0보다 큰 지갑을 $\log_{10}$ 눈금으로 그렸고, 지갑 수도 로그 눈금이에요. 0점 지갑은 뺐어요. 이 지갑들에는 그 그래프의 화살표가 하나도 닿지 않아요. 엔도스랭크는 $966$개, AWP는 $381$개예요. 엔도스랭크에서는 $4{,}424$개 지갑이 허락을 받지 않은 지갑의 점수 $9.62\times10^{-5}$를 함께 가지고, 그보다 위에 있는 것은 $131$개뿐이에요.}
\label{fig:score-distributions}
\end{figure}

\begin{figure}[htbp]
\centering
\includegraphics[width=0.8\textwidth]{\figpath{rank-scatter.pdf}}
\caption[{매칭 코호트에서 엔도스랭크와 AWP의 순위--순위 산점도예요(지갑 $2{,}000$개를 무작위로 뽑았어요). 점선 대각선(두 순위가 같은 선) 위의 점은 두 순서가 똑같다는 뜻이에요.}]{매칭 코호트에서 엔도스랭크와 AWP를 맞대어 그린 순위--순위 산점도예요(지갑 $2{,}000$개를 무작위로 뽑았어요). 동점인 점수는 평균 순위를 함께 가져요. 켄달의 $\tau_b$도 동점을 이렇게 다뤄요. 세로 기둥 두 개는 엔도스랭크의 동점 덩어리예요. 하나는 공통 점수 $9.62\times10^{-5}$를 가진 $4{,}424$개 지갑이고, 다른 하나는 0점인 $966$개 지갑이에요. 점선 대각선 위의 점은 두 순서에서 순위가 같아요.}
\label{fig:rank-scatter}
\end{figure}

H1은 뒷받침돼요. 두 점수는 코호트를 다르게 줄 세우고, 두 점수가 관련되는 것은 주로 엔도스랭크가 0점을 주는 지갑들을 통해서예요. 각 점수가 동점 말고 무엇을 담고 있는지는 기간 밖에서 드러나요. 기간 밖에서는 모든 스펜더가 허락을 받았고, 엔도스랭크에 동점이 거의 남지 않아요(\ref{sec:holdout-results}절). 송금에 바탕을 둔 점수에 기대는 서비스는, 허락에 바탕을 둔 점수로 바꾸거나 그런 점수와 섞어도 지갑 순서가 그대로 남는다고 생각할 수 없어요.

\dissertationSubheading[sec:robustness-results]{강건성 점검과 민감도 점검}

세 가지 점검은 결과가 감쇠 값(화살표를 따라갈 확률) 하나, 토큰 구성 하나, 표본 정의 하나에 달려 있는지를 물어요. 이 점검들은 엔도스랭크와 AWP에 대해 했어요. 처음 두 점검은 표본을 그대로 두고 설정값이나 토큰 묶음을 조금 바꿔요. 세 번째 점검은 표본 정의 자체를 바꾸는데, 여기서는 계수가 움직일 것으로 예상해요. 네 번째 점검은 엔도스랭크가 재시작 규칙에 얼마나 기대는지를 물어요.

\dissertationSubsubheading[sub:damping]{감쇠}

% Real BigQuery data -- regenerate with:
%   2-Dissertation-Draft-Works/wallet-reputation-experiments/scripts/run_dissertation_eval.py --real --export-latex
% Generated by export_latex_results.py -- do not edit by hand unless re-export fails
\begin{table}[htbp]
\centering
\caption[Robustness: mean Kendall $\tau$ under PageRank damping $d \in \{0.75, 0.85, 0.95\}$ on the matched alignment cohort.]{Robustness: mean Kendall $\tau$ under PageRank damping $d \in \{0.75, 0.85, 0.95\}$ on the matched alignment cohort ($n=5{,}521$). The $d=0.85$ runtimes are the same measurements as Table~\ref{tab:benchmark-runtime} (one timing session; each graph--damping pair is timed once).}
\label{tab:robustness-damping}
\fitwidth{%
\begin{tabular}{rcccccc}
\toprule
$d$ & ER all. & AWP all. & ER tr. & AWP tr. & ER (s) & AWP (s) \\
\midrule
0.75 & 0.375 & -0.022 & 0.330 & 0.626 & 0.541 & 5.011 \\
0.85 & 0.375 & -0.024 & 0.330 & 0.612 & 0.547 & 5.029 \\
0.95 & 0.375 & -0.030 & 0.330 & 0.585 & 0.541 & 5.187 \\
\bottomrule
\end{tabular}
}
\end{table}


엔도스랭크는 $d\in\{0.75,0.85,0.95\}$에 걸쳐 가족 평균이 소수 셋째 자리에서도 움직이지 않아요. AWP는 $d$가 커지면서 송금 계수가 $0.626$에서 $0.585$로 떨어지고, 허락 계수는 조금 음수인 채로 머물러요. 두 점수의 모습 모두, 흔히 쓰는 $d=0.85$를 골라서 생긴 것이 아니에요. 엔도스랭크가 흔들리지 않는 것은 대부분 구조 때문이에요. 엔도스랭크의 공통 값보다 위에 있는 코호트 지갑은 $131$개뿐이에요(\ref{sec:rank-divergence}절). $d$를 바꾸면 공통 값은 움직이지만, 그 값을 함께 가진 $4{,}424$개 지갑은 여전히 동점이에요. 그래서 순위 상관(두 순위가 얼마나 비슷한지)은 거의 움직일 수 없어요. 실행 시간도 거의 움직이지 않아요. 엔도스랭크는 $d=0.75$, $0.85$, $0.95$에서 $0.541$, $0.547$, $0.541$\,s가 걸리고, AWP는 $5.011$, $5.029$, $5.187$\,s가 걸려요. 그래서 비율은 $9.2\times$에서 $9.6\times$ 사이에 머물러요. 그래프와 감쇠 값의 짝마다 한 세션에서 한 번씩 시간을 쟀어요. 호출마다 대부분의 시간은 행렬 조립이 차지하고, 조립은 $d$와 상관없어요. 그래서 $d=0.95$에서 늘어난 반복이 시간을 아주 조금만 늘려요. 그림~\ref{fig:damping-sensitivity}는 가족 평균을 그려요.

\begin{figure}[htbp]
\centering
\includegraphics[width=0.85\textwidth]{\figpath{damping-sensitivity.pdf}}
\caption[{$d\in\{0.75,0.85,0.95\}$에서 허락 가족과 송금 가족의 가족 평균 켄달 $\tau$의 민감도(표~\ref{tab:robustness-damping}의 값).}]{$d\in\{0.75,0.85,0.95\}$에 걸친, 허락 가족과 송금 가족의 가족 평균 켄달 $\tau$의 감쇠 민감도(표~\ref{tab:robustness-damping}의 값).}
\label{fig:damping-sensitivity}
\end{figure}

\dissertationSubsubheading[sub:token-subgraph]{상위 토큰 부분 그래프}

% Real BigQuery data -- regenerate with:
%   2-Dissertation-Draft-Works/wallet-reputation-experiments/scripts/run_dissertation_eval.py --real --export-latex
% Generated by export_latex_results.py -- do not edit by hand unless re-export fails
\begin{table}[htbp]
\centering
\caption{Robustness: mean Kendall $\tau$ on the matched cohort when both graphs keep only 20 ERC-20 tokens, chosen by summed raw amount or by number of rows (transfer events and latest allowances), next to the value on all tokens from Table~\ref{tab:alignment-hybrid}.}
\label{tab:robustness-tokens}
\fitwidth{%
\begin{tabular}{lccccccccc}
\toprule
Method & \multicolumn{3}{c}{Transfer} & \multicolumn{3}{c}{Allowance} & \multicolumn{3}{c}{Sybil} \\
\cmidrule(lr){2-4} \cmidrule(lr){5-7} \cmidrule(lr){8-10}
 & by amount & by count & all & by amount & by count & all & by amount & by count & all \\
\midrule
EndorseRank & 0.328 & 0.329 & 0.330 & 0.370 & 0.370 & 0.375 & 0.231 & 0.230 & 0.231 \\
AWP & 0.598 & 0.603 & 0.612 & -0.021 & -0.022 & -0.024 & 0.273 & 0.275 & 0.278 \\
\bottomrule
\end{tabular}
}
\end{table}


두 그래프를 토큰 스무 개로 좁히면(표~\ref{tab:robustness-tokens}), 가족 평균이 가장 크게 바뀐 것은 $0.014$예요. 원래 단위 금액의 합이 가장 큰 스무 개 토큰을 골랐을 때, 송금 가족에서 AWP가 바뀐 크기예요. 기록 줄이 가장 많은 스무 개 토큰을 고르면 가장 큰 변화는 $0.009$예요. 엔도스랭크는 많아야 $0.005$ 움직여요. 두 목록에는 같은 토큰이 열한 개 있어요. 그래서 이 모습은 두 목록 어디에도 들지 않은 수많은 토큰에 달려 있지 않아요. 두 목록 모두 경제적 가치로 토큰을 줄 세운 것이 아니에요. 원래 단위 금액의 합은 무제한 허락(얼마든지 써도 된다는 허락)이 많은 토큰에 유리하고, 줄 수는 주인이 자주 바뀌는 토큰에 유리해요.

\dissertationSubsubheading[sub:sample-size]{표본 정의 민감도}

% Real BigQuery data -- regenerate with:
%   2-Dissertation-Draft-Works/wallet-reputation-experiments/scripts/run_dissertation_eval.py --real --export-latex
% Generated by export_latex_results.py -- do not edit by hand unless re-export fails
\begin{table}[htbp]
\centering
\caption[Sensitivity of alignment to the evaluation set: mean Kendall $\tau$ and runtime at staged subsample sizes drawn from the expanded wallet pool.]{Sensitivity of alignment to the evaluation set: mean Kendall $\tau$ and runtime at staged subsample sizes drawn from the expanded wallet pool ($N=99{,}080$; seed \texttt{benchmark-tier2-v1}). Rows are evaluated on different wallet sets from the matched cohort, so $\tau$ values are not comparable with Table~\ref{tab:alignment-hybrid}; they show how much the statistic moves with the evaluation population. Runtimes are the same measurements as Table~\ref{tab:benchmark-scaling}.}
\label{tab:robustness-sample-size}
\fitwidth{%
\begin{tabular}{rccccrrc}
\toprule
$n$ & ER all. & AWP all. & ER tr. & AWP tr. & ER (s) & AWP (s) & Speedup \\
\midrule
10{,}000 & 0.343 & 0.239 & 0.821 & 0.900 & 0.075 & 1.078 & 14.37 \\
50{,}000 & 0.333 & 0.230 & 0.827 & 0.909 & 0.342 & 3.414 & 9.98 \\
99{,}080 & 0.332 & 0.230 & 0.826 & 0.909 & 0.547 & 5.029 & 9.19 \\
\bottomrule
\end{tabular}
}
\end{table}


표~\ref{tab:robustness-sample-size}에서는 실행 시간 규모 시험에 쓴 단계별 풀 부분 표본에서 송금 가족과 허락 가족의 평균을 다시 계산해요. 송금 가족 평균은 매칭 코호트에서보다 훨씬 크게 나오고($0.82$--$0.91$), AWP의 허락 계수는 양수로 바뀌어요($0.23$--$0.24$). 반면 엔도스랭크의 허락 계수($0.332$--$0.343$)는 코호트 값 $0.375$보다 낮게 끝나요. 이 변화들은 어느 것도 점수에 대한 증거가 아니에요. 풀 부분 표본에는 한 그래프나 두 그래프 모두에 화살표가 없는 지갑이 들어와요. 점수와 지표를 지갑별로 맞춰 놓고 나면 그런 지갑은 점수도 0이고 대리 지표도 0이에요. 그리고 두 순위의 바닥에 몰린 0의 동점들은 켄달 $\tau$를 기계적으로 끌어올려요(\ref{sub:sample-size-protocol}절). 그래서 이 표는 표본 정의의 민감도 분석으로 보고해요. 이 표의 값을 표~\ref{tab:alignment-hybrid}의 값과 견주면 안 돼요. 모든 풀 단계에서 엔도스랭크 그래프의 점 수는 표본의 지갑 수보다 훨씬 적어요(표~\ref{tab:benchmark-scaling}). 그래서 풀 지갑 대부분이 엔도스랭크에서 0점을 받아요. 코호트에서는 $5{,}521$개 가운데 $966$개가 0점이에요.

감쇠 점검과 상위 토큰 점검을 보면, AWP의 같은 기간 계수는 설정값이나 토큰 구성에 달려 있지 않아요. 엔도스랭크의 계수는 거의 움직일 수 없었어요. 어떤 $d$나 토큰 묶음을 써도 동점이 그대로 남기 때문이에요. 표본 정의 점검처럼 대상 집단이 바뀌면, 값의 크기 그 자체(절댓값)는 그대로 남지 않아요.

\dissertationSubsubheading[sub:endorserank-restarts]{엔도스랭크의 재시작}

엔도스랭크는 걷기를 모든 점에서 고르게 다시 출발시키지만, AWP는 주소들이 최근에 얼마나 많이 보냈는지에 비례해서 그 주소에서 다시 출발시켜요(\ref{sub:endorserank-graph}절). AWP의 규칙을 허락에 옮겨 오면, 주인들이 최근에 허락해 준 스펜더 수에 비례해서 그 주인에서 다시 출발하게 돼요. 표~\ref{tab:endorserank-restarts}는 두 가지 엔도스랭크의 점수를 매기고, 그 옆에 AWP를 함께 놓아요. 재시작 규칙은 이 비교를 한 뒤에 골랐어요.

% Real BigQuery data -- regenerate with:
%   2-Dissertation-Draft-Works/wallet-reputation-experiments/scripts/run_dissertation_eval.py --real --export-latex
% Generated by export_latex_results.py -- do not edit by hand unless re-export fails
\begin{table}[htbp]
\centering
\caption[EndorseRank with uniform restarts and with restarts weighted by approving activity, AWP's rule applied to the allowance graph, next to AWP: family-mean Kendall $\tau$ on the matched cohort ($n=5{,}521$) and $\tau$ with the two spring labels on the spender cohort.]{EndorseRank with uniform restarts and with restarts weighted by approving activity, AWP's rule applied to the allowance graph, next to AWP: family-mean Kendall $\tau$ on the matched cohort ($n=5{,}521$) and $\tau$ with the two spring labels on the spender cohort ($n=1{,}335$). Both EndorseRank rows weight each latest allowance by $\sigma(\Delta t)V(z)$. CI = 95\% percentile bootstrap over wallets (400 paired resamples, seed 42).}
\label{tab:endorserank-restarts}
\fitwidth{%
\begin{tabular}{llccccc}
\toprule
 & & \multicolumn{3}{c}{Same window, matched cohort} & \multicolumn{2}{c}{Spring holdout, spenders} \\
\cmidrule(lr){3-5} \cmidrule(lr){6-7}
Method & Restarts & Transfer & Allowance & Sybil & New approval pairs & New senders \\
\midrule
EndorseRank & uniform & \begin{tabular}[c]{@{}c@{}}0.330\\[-1pt]{\footnotesize [0.311, 0.350]}\end{tabular} & \begin{tabular}[c]{@{}c@{}}0.375\\[-1pt]{\footnotesize [0.342, 0.401]}\end{tabular} & \begin{tabular}[c]{@{}c@{}}0.231\\[-1pt]{\footnotesize [0.212, 0.250]}\end{tabular} & \begin{tabular}[c]{@{}c@{}}0.394\\[-1pt]{\footnotesize [0.360, 0.431]}\end{tabular} & \begin{tabular}[c]{@{}c@{}}0.283\\[-1pt]{\footnotesize [0.247, 0.320]}\end{tabular} \\
EndorseRank & by approving activity & \begin{tabular}[c]{@{}c@{}}0.445\\[-1pt]{\footnotesize [0.433, 0.457]}\end{tabular} & \begin{tabular}[c]{@{}c@{}}0.044\\[-1pt]{\footnotesize [0.029, 0.062]}\end{tabular} & \begin{tabular}[c]{@{}c@{}}0.245\\[-1pt]{\footnotesize [0.233, 0.259]}\end{tabular} & \begin{tabular}[c]{@{}c@{}}0.372\\[-1pt]{\footnotesize [0.327, 0.417]}\end{tabular} & \begin{tabular}[c]{@{}c@{}}0.301\\[-1pt]{\footnotesize [0.254, 0.348]}\end{tabular} \\
AWP & by sending activity & \begin{tabular}[c]{@{}c@{}}0.612\\[-1pt]{\footnotesize [0.602, 0.622]}\end{tabular} & \begin{tabular}[c]{@{}c@{}}-0.024\\[-1pt]{\footnotesize [-0.043, -0.002]}\end{tabular} & \begin{tabular}[c]{@{}c@{}}0.278\\[-1pt]{\footnotesize [0.264, 0.292]}\end{tabular} & \begin{tabular}[c]{@{}c@{}}0.123\\[-1pt]{\footnotesize [0.068, 0.177]}\end{tabular} & \begin{tabular}[c]{@{}c@{}}0.344\\[-1pt]{\footnotesize [0.307, 0.383]}\end{tabular} \\
\bottomrule
\end{tabular}
}
\end{table}


AWP의 재시작을 쓰면, 엔도스랭크의 같은 기간 허락 계수는 $0.375$에서 $0.044$로 떨어지고, 송금 계수는 $0.330$에서 $0.445$로 올라가요. 이때 재시작 몫은 허락을 가장 많이 해 주는 주인들에게 가요. 코호트의 트레이더들이 바로 그런 주인이에요. 그래서 이 점수는 트레이더들을 얼마나 많이 허락해 주는지로 줄 세우는데, 이것은 그들이 받는 보증보다 그들의 활동에 더 가까워요. 기간 밖에서는 두 가지가 비슷해요. 봄(3월부터 5월까지) 스펜더들에서 두 가지는 새 허락 쌍에서 $0.394$와 $0.372$, 새 송금자(새로 송금을 보내온 주소)에서 $0.283$과 $0.301$이고, 구간이 서로 겹쳐요. 6월부터 8월까지는 $0.333$과 $0.325$, 그리고 $0.255$와 $0.284$예요(표~\ref{tab:fresh-posthoc}). 고르게 다시 출발하는 규칙은 같은 기간에서 허락 구성개념을 지키면서, 기간 밖에서는 잴 수 있을 만큼 잃는 것이 없어요.

\dissertationSubheading[sec:holdout-results]{시간 홀드아웃(RQ4, RQ5)}

같은 기간 점검으로는 점수를 그 점수 자신의 구성개념과 떼어 놓을 수 없어요. 그래서 홀드아웃은 그 대신, 어떤 날까지 잰 중심성(그래프에서 얼마나 중요한 자리에 있는지)이 그날 뒤의 행동을 미리 맞히는지를 시험해요. \ref{sec:holdout-protocol}절을 따라 스펜더 코호트의 점수는 2026년 2월 28일에 고정하고, 그다음 2026년 3월 1일부터 5월 31일까지 지갑마다 라벨(나중에 채점에 쓰는 정답 값) 두 개를 세요. 하나는 새 주인--스펜더 허락 쌍(보증 구성개념)이고, 다른 하나는 새 송금자(흐름 구성개념)예요. 표~\ref{tab:holdout-spenders}은 $t_1$에 계산한 모든 점수를, 원시 차수 기준선(비교 기준) 두 개까지 포함해서, 두 라벨에 견주어요. 표~\ref{tab:holdout-diff}은 짝지은 대비를 보여 줘요. 먼저 결합 점수를 계산하기 전에 정한 대비가 나오고, 다음으로 라벨을 본 뒤에 계산한 대비가 나와요.

% Real local parquet -- regenerate with:
%   2-Dissertation-Draft-Works/wallet-reputation-experiments/scripts/run_holdout_eval.py --cohort spenders
% Generated by export_latex_results.py -- do not edit by hand unless re-export fails
\begin{table}[htbp]
\centering
\caption[Temporal holdout on the spender cohort ($n=1{,}335$): Kendall $\tau$ between scores frozen at 28~February~2026 and two labels counted on the wallet between 1~March~2026 and 31~May~2026: the number of new owner--spender approval pairs and the number of addresses that sent a transfer to the wallet for the first time.]{Temporal holdout on the spender cohort ($n=1{,}335$): Kendall $\tau$ between scores frozen at 28~February~2026 and two labels counted on the wallet between 1~March~2026 and 31~May~2026: the number of new owner--spender approval pairs and the number of addresses that sent a transfer to the wallet for the first time. CI = 95\% percentile bootstrap over wallets (400 paired resamples). 222 wallets received at least one new pair and 175 at least one new sender. Baseline rows are the raw degrees at $t_1$ that each single-layer PageRank smooths. C-PR ($\lambda=1$) and C-PR ($\lambda=0$) are the allowance and transfer layers of the coupled operator walked alone.}
\label{tab:holdout-spenders}
\fitwidth{%
\begin{tabular}{lcc}
\toprule
Score at $t_1$ & New approval pairs ($\tau$ [95\% CI]) & New transfer senders ($\tau$ [95\% CI]) \\
\midrule
In-approve degree at $t_1$ (baseline) & 0.711 [0.665, 0.759] & 0.404 [0.348, 0.467] \\
Transfer in-degree at $t_1$ (baseline) & 0.117 [0.049, 0.191] & 0.472 [0.423, 0.520] \\
\midrule
EndorseRank & 0.394 [0.360, 0.431] & 0.283 [0.247, 0.320] \\
AWP & 0.123 [0.068, 0.177] & 0.344 [0.307, 0.383] \\
\midrule
C-PR ($\lambda=1$) & 0.479 [0.427, 0.529] & 0.372 [0.325, 0.424] \\
C-PR ($\lambda=0$) & 0.044 [-0.011, 0.098] & 0.270 [0.235, 0.310] \\
C-PR ($\lambda=0.25$) & 0.271 [0.231, 0.314] & 0.266 [0.231, 0.306] \\
C-PR ($\lambda=0.5$) & 0.295 [0.252, 0.336] & 0.277 [0.241, 0.318] \\
C-PR ($\lambda=0.75$) & 0.314 [0.271, 0.357] & 0.287 [0.253, 0.330] \\
S-PR & 0.114 [0.052, 0.164] & 0.264 [0.229, 0.304] \\
\bottomrule
\end{tabular}
}
\end{table}


% Real local parquet -- regenerate with:
%   2-Dissertation-Draft-Works/wallet-reputation-experiments/scripts/run_holdout_eval.py --cohort spenders
% Generated by export_latex_results.py -- do not edit by hand unless re-export fails
\begin{table}[htbp]
\centering
\caption[Paired bootstrap differences on the spender holdout.]{Paired bootstrap differences on the spender holdout ($n=1{,}335$, 400 resamples shared by $a$ and $b$). The first block was fixed before the run: the increment of C-PR ($\lambda=1$), the allowance layer walked alone, over its raw degree at $t_1$, and the hybrid contrasts fixed on 14~September~2026, which compare C-PR and S-PR with the two layers they are built from, C-PR ($\lambda=1$) and C-PR ($\lambda=0$). The second block was computed after the labels were known: the increments of EndorseRank and AWP over their raw degrees, C-PR against C-PR ($\lambda=0$) on new approval pairs, and C-PR against EndorseRank and AWP. Primary criterion of 14~September (C-PR no worse than C-PR ($\lambda=1$) on new approval pairs and better than C-PR ($\lambda=0$) on new transfer senders): not met. The differences between EndorseRank and AWP are in Table~\ref{tab:endorserank-awp-holdout}.}
\label{tab:holdout-diff}
\fitwidth{%
\begin{tabular}{llccccc}
\toprule
Contrast ($a$ minus $b$) & Label & $\tau_a$ & $\tau_b$ & $\Delta\tau$ & 95\% CI of $\Delta\tau$ & $P(\Delta\tau>0)$ \\
\midrule
\multicolumn{7}{l}{\emph{Fixed before the run}} \\
C-PR ($\lambda=1$) minus t1 in-approve degree & new approval pairs & 0.479 & 0.711 & -0.233 & [-0.291, -0.183] & 0.000 \\
C-PR minus C-PR ($\lambda=1$) & new approval pairs & 0.295 & 0.479 & -0.184 & [-0.213, -0.149] & 0.000 \\
C-PR minus C-PR ($\lambda=0$) & new transfer senders & 0.277 & 0.270 & +0.007 & [-0.012, 0.028] & 0.735 \\
S-PR minus C-PR ($\lambda=1$) & new approval pairs & 0.114 & 0.479 & -0.365 & [-0.431, -0.300] & 0.000 \\
S-PR minus C-PR ($\lambda=0$) & new transfer senders & 0.264 & 0.270 & -0.006 & [-0.017, 0.006] & 0.163 \\
\midrule
\multicolumn{7}{l}{\emph{Computed after the labels were known}} \\
EndorseRank minus t1 in-approve degree & new approval pairs & 0.394 & 0.711 & -0.317 & [-0.354, -0.283] & 0.000 \\
AWP minus t1 in-degree & new transfer senders & 0.344 & 0.472 & -0.127 & [-0.150, -0.099] & 0.000 \\
C-PR minus C-PR ($\lambda=0$) & new approval pairs & 0.295 & 0.044 & +0.251 & [0.203, 0.306] & 1.000 \\
C-PR minus EndorseRank & new approval pairs & 0.295 & 0.394 & -0.099 & [-0.133, -0.064] & 0.000 \\
C-PR minus AWP & new transfer senders & 0.277 & 0.344 & -0.067 & [-0.095, -0.039] & 0.000 \\
C-PR minus AWP & new approval pairs & 0.295 & 0.123 & +0.172 & [0.111, 0.226] & 1.000 \\
\bottomrule
\end{tabular}
}
\end{table}


두 가지 결과를 차례로 살펴봐요.

첫 번째 결과는 RQ4의 두 번째 부분에 “아니요”라고 답해요. 어떤 페이지랭크도 $t_1$에서 자기 화살표 종류의 원시 차수보다 낫지 않아요. $t_1$의 받은 허락 수는 새 허락 쌍을 $\tau=0.711$로 미리 맞혀요. 허락 그래프를 걷는 점수들이 가장 가깝지만, 그래도 한참 못 미쳐요. C-PR ($\lambda=1$)은 $0.479$로, 짝지은 차이가 $-0.233$이에요(구간 $[-0.291,-0.183]$, 돌리기 전에 미리 정한 대비). 엔도스랭크는 $0.394$로, 차이가 $-0.317$이에요($[-0.354,-0.283]$). $t_1$의 송금 입차수는 새 송금자를 $\tau=0.472$로 미리 맞혀요. AWP는 $0.344$이고($\Delta\tau=-0.127$, 구간 $[-0.150,-0.099]$), C-PR ($\lambda=0$)은 $0.270$이에요. 모든 구간이 0보다 한참 아래에 있어요. 석 달 기간 안에서는, 이미 거래 상대가 많은 지갑일수록 새 거래 상대도 더 많이 얻어요. 그리고 그 자리에서 센 수가 이것을 바로 잡아내요. 페이지랭크는 이 수를, 무게를 달고 퍼뜨린 양으로 바꿔요. 엔도스랭크에서는 한 주인의 보증이 그 주인이 허락한 모든 스펜더에게 나뉘고, AWP에서는 한 송금자의 몫이 그 송금자가 돈을 보낸 모든 주소에게 나뉘어요. 이 라벨들에서는 그렇게 나누는 일이 잡음처럼 작용해요. 페이지랭크가 차수 말고 무언가를 더한다고 해도, 그것은 자기 구성개념이 앞으로 자라는 것을 미리 맞히는 데에서는 드러나지 않아요.

표를 가로질러 읽으면, 각 라벨을 가장 잘 줄 세우는 것은 그 라벨과 같은 화살표 종류의 원시 차수예요. 새 허락 쌍에서는 허락 그래프를 걷는 두 점수가 그다음이고, 그다음이 중간 비율로 섞은 점수들($0.271$에서 $0.314$까지)이에요. 반면 송금 입차수, AWP, C-PR ($\lambda=0$), S-PR은 $0.123$ 이하에 머물러요. 새 송금자에서는 순서가 덜 깔끔해요. 받은 허락 수($0.404$)와 C-PR ($\lambda=1$)($0.372$)이 송금을 걷는 어떤 점수보다도 이 라벨을 잘 줄 세워요. 송금을 걷는 점수 가운데에서는 AWP가 가장 높아요($0.344$). 이 코호트에서는 많은 주인이 허락해 준 스펜더가 새 송금자도 많이 얻어요. 여기서 엔도스랭크는 동점을 거의 남기지 않아요. 스펜더 $1{,}335$개에 서로 다른 점수를 $1{,}291$개 주고, 가장 큰 동점 덩어리에는 $42$개가 들어 있어요. 그래서 엔도스랭크의 계수는 거의 모든 스펜더의 순서에 바탕을 둬요. 허락 층을 금액 그대로 걷는 C-PR ($\lambda=1$)은 두 라벨 모두 엔도스랭크보다 높게 줄 세워요. 그 줄들은 무제한 허락이 정해요(\ref{sub:coupled-graph}절). 그래서 이 코호트에서는, 허락을 하나마다 한 번씩만 세면 놓치는 신호의 일부를 금액 그대로가 지켜 냈어요. 왜 그런지는 이 논문에서 시험하지 않아요.

걷는 점수들이 차수에 못 미치는 까닭은 포함 범위(점수가 닿는 지갑의 범위)로 설명되지 않아요. 스펜더 1,335개 가운데 318개는 동결일(점수를 고정한 날)까지 송금 화살표가 하나도 없어서 AWP에서 0점이에요. 이 스펜더들은 나머지보다 새 허락 쌍을 더 자주 얻었고(22\% 대 15\%), 새 송금자는 덜 자주 얻었어요(4\% 대 16\%). 그래서 AWP의 0점들은 첫 번째 라벨에서는 AWP의 계수를 낮추고, 두 번째 라벨에서는 높여요. 표~\ref{tab:holdout-receiving}는 홀드아웃을, 동결일 전에 다른 주소에게서 송금을 받은 적이 있는 스펜더 977개로 좁혀요. 여기서도 AWP는 송금 입차수보다 한참 낮고($0.388$ 대 $0.701$; $-0.314$, $[-0.353,-0.276]$), 엔도스랭크는 받은 허락 수보다 낮아요($0.340$ 대 $0.679$). 이렇게 좁히는 방법은 결과를 안 뒤에 골랐어요.

% Post hoc checks -- regenerate with:
%   2-Dissertation-Draft-Works/wallet-reputation-experiments/scripts/run_supplementary_checks.py
% Generated by export_latex_results.py -- do not edit by hand unless re-export fails
\begin{table}[htbp]
\centering
\caption[스펜더 1{,}335개 가운데, 동결일 전에 다른 주소에서 송금을 받은 적이 있는 977개로 좁힌 스펜더 홀드아웃.]{스펜더 1{,}335개 가운데, 동결일 전에 다른 주소에서 송금을 받은 적이 있는 977개로 좁힌 스펜더 홀드아웃이에요. 나머지 358개에는 들어오는 송금 화살표가 없고, 스펜더 318개는 AWP에서 점수가 0이에요. 이렇게 좁히는 것은 결과를 본 뒤에 해 보았어요. 이 묶음에서 엔도스랭크와 AWP 사이의 차이는 표~\ref{tab:endorserank-awp-holdout}에 있어요. 신뢰구간은 지갑을 다시 뽑아 구한 95\% 백분위 부트스트랩 구간이에요(다시 뽑기 400번, 시드 42).}
\label{tab:holdout-receiving}
\fitwidth{%
\begin{tabular}{lcc}
\toprule
$t_1$의 점수 & 새 허락 쌍($\tau$ [95\% 신뢰구간]) & 새 송금자($\tau$ [95\% 신뢰구간]) \\
\midrule
EndorseRank & 0.340 [0.299, 0.386] & 0.358 [0.316, 0.396] \\
AWP & 0.298 [0.252, 0.343] & 0.388 [0.348, 0.430] \\
\midrule
$t_1$의 받은 허락 수 & 0.679 [0.616, 0.736] & 0.570 [0.505, 0.635] \\
$t_1$의 송금 입차수 & 0.569 [0.510, 0.631] & 0.701 [0.651, 0.750] \\
\midrule
AWP 빼기 t1 입차수(새 송금자) & \multicolumn{2}{c}{$\Delta\tau=-0.314$ [-0.353, -0.276]} \\
\bottomrule
\end{tabular}
}
\end{table}


두 번째 결과는, 결합 연산자가 9월 14일의 규칙을 통과하지 못한다는 거예요. 다만 탐색적(미리 정하지 않고 살펴본) 비교에서는 송금 층만 걷는 것보다 나아요. 그 규칙은 C-PR이 새 허락 쌍에서는 C-PR ($\lambda=1$)만큼 하고, 새 송금자에서는 C-PR ($\lambda=0$)을 이기기를 요구했어요. $\lambda=0.5$에서 C-PR은 새 허락 쌍에서 허락 걷기보다 못하고($\Delta\tau=-0.184$, 구간 $[-0.213,-0.149]$), 새 송금자에서는 송금 걷기와 다르지 않아요($\Delta\tau=+0.007$, 구간 $[-0.012,0.028]$). 그래서 규칙의 주된 부분에 있는 두 조건이 모두 통과하지 못해요. 규칙에 들어 있지 않던 비교, 곧 송금 걷기 하나와만 견주면 C-PR이 더 나아요. 새 허락 쌍에서 C-PR은 $0.295$, 송금 걷기는 $0.044$예요. 짝지은 차이는 $+0.251$이에요(구간 $[0.203,0.306]$, 같은 다시 뽑은 표본으로 구했어요). 그리고 새 송금자에서는 알아챌 만한 손해가 없어요. 이 대비는 라벨을 본 뒤에 계산했어요. 그래서 가설로 다루고, 아래의 등록한 기간에서 다시 시험해요. 엔도스랭크와 AWP와의 대비도 라벨을 본 뒤에 계산했어요. 새 허락 쌍에서 C-PR은 두 점수 사이에 있어요. 엔도스랭크보다 $0.099$ 낮고($[-0.133,-0.064]$), AWP보다 $0.172$ 높아요($[0.111,0.226]$). 그리고 새 송금자에서는 AWP에 뒤져요($-0.067$, $[-0.095,-0.039]$). 그러니까 봄 자료에서 F2(등록한 규칙의 두 번째 가설)와 같은 시험을 AWP와 견주어 하면 통과하지 못했을 거예요. 등록 재현은 그 비교를 결정과 나란히 보고해요(\ref{sub:fresh-holdout}절). $\lambda$의 세 값에 걸쳐 두 라벨 모두 순서가 한 방향으로만 움직여요. 새 허락 쌍에서 $\tau$는 $0.271$($\lambda=0.25$)에서 $0.295$를 거쳐 $0.314$($\lambda=0.75$)로 올라가고, 새 송금자에서는 $0.266$에서 $0.277$을 거쳐 $0.287$로 올라가요. 허락 쪽 무게를 늘리면 두 라벨 모두에서 도움이 돼요. 이 코호트에서는 송금 층이, 허락 층이 담은 것에 기간 밖 신호를 하나도 더하지 않아요. 씨앗을 쓴 빼 보기 실험은 허락 신호가 어디에 있어야 하는지 보여 줘요. S-PR은 흐름 라벨에서 송금 걷기와 같은 정도이고($\Delta\tau=-0.006$, 구간 $[-0.017,0.006]$), 보증 라벨에서는 허락 걷기보다 한참 낮아요($\Delta\tau=-0.365$). 송금 화살표가 걷기가 갈 곳을 정할 때에는, 보증에서 끌어낸 재시작 분포가 거의 아무것도 바꾸지 못해요. 스펜더 홀드아웃에서 허락 신호는 허락 화살표가 걷기를 실어 나를 때 살아남아요. 허락 걷기에서는 온전히, C-PR에서는 묽어진 채로 살아남아요. 그리고 허락 화살표가 걷기가 다시 출발하는 곳만 정할 때에는 사라져요.

표 전체에 주의할 점이 두 가지 있어요. 각 라벨은 같은 방식으로 만든 것을 미래로 옮긴 거예요(새 허락, 새 송금자). 그래서 이 결과들은 어느 것도 신용이나 손실과는 관계가 없어요. 또 두 라벨 모두 듬성듬성해요. 지갑 $1{,}335$개 가운데 새 허락 쌍을 얻은 것은 $222$개, 새 송금자를 얻은 것은 $175$개예요. 모든 계수는 오로지 이 작은 묶음이 어떤 순서로 서는지에 달려 있어요.

홀드아웃에서는 다른 후보 라벨들도 계산했어요. 여기서 보고하지만, 뒷받침하는 증거로 내놓지는 않아요. 허락 취소 횟수와 빼내기 어림 규칙(허락 뒤에 토큰이 빠져나간 일을 찾는 대략의 규칙)은 둘 다 엔도스랭크와 함께 올라가요($\tau=0.269$, 구간 $[0.226,0.308]$, 그리고 $0.422$, 구간 $[0.389,0.455]$). 받은 허락 수 그대로와는 더 가파르게 올라가요($0.483$과 $0.575$). 보증으로 읽으려면 부호가 틀렸어요. 믿을 만한 스펜더라면 허락을 더 자주가 아니라 덜 자주 잃어야 하기 때문이에요. 더 그럴듯한 해석은 양(활동량)이에요. 주인이 많은 스펜더는 주인들이 하는 모든 일을 더 많이 겪는데, 허락 취소와 빼내기도 마찬가지예요. 그래서 이 라벨들로는 믿음과 쓰임을 갈라낼 수 없어요. 이 스펜더 묶음에서는 Aave(대출 서비스) 청산이 하나도 없어요.

\dissertationSubsubheading[sub:trader-results]{트레이더 결과(탐색적)}

\ref{sub:trader-holdout}절의 트레이더 실행은 라벨을 본 뒤에 점수를 매겼으므로, 탐색적 결과로 보고해요. 표~\ref{tab:holdout-traders}은 3월부터 5월까지 포지션을 적어도 세 번 닫은 매칭 지갑 1,860개의 GMX 라벨에서 모든 점수를 보여 줘요.

% Real local parquet -- regenerate with:
%   2-Dissertation-Draft-Works/wallet-reputation-experiments/scripts/run_holdout_eval.py --cohort matched
% Generated by export_latex_results.py -- do not edit by hand unless re-export fails
\begin{table}[htbp]
\centering
\caption[봄 트레이더 실행: 2026년 2월 28일에 고정한 점수와 2026년 3월 1일부터 2026년 5월 31일까지 센 GMX~V2 라벨 사이의 켄달 $\tau$(그 기간에 닫기를 세 번 이상 한 매칭 지갑 1{,}860개).]{봄 트레이더 실행의 결과예요. 2026년 2월 28일에 고정한 점수와, 2026년 3월 1일부터 2026년 5월 31일까지 센 GMX~V2 라벨 사이의 켄달 $\tau$를, 그 기간에 닫기를 세 번 이상 한 매칭 지갑 1{,}860개에서 구했어요. 이 실행은 라벨을 본 뒤에 점수를 매겼어요. 이익을 낸 닫기 수와 실현 이익은 닫기 수가 많을수록 커지지만, 두 비율은 그렇지 않아요. 신뢰구간은 지갑을 다시 뽑아 구한 95\% 백분위 부트스트랩 구간이에요(짝지은 다시 뽑기 400번, 시드 42).}
\label{tab:holdout-traders}
\fitwidth{%
\begin{tabular}{lcccc}
\toprule
$t_1$의 점수 & 이익을 낸 닫기 수 & 실현 이익 & 이익을 낸 비율 & 손실 없는 비율 \\
\midrule
$t_1$의 받은 허락 수(기준선) & \begin{tabular}[c]{@{}c@{}}-0.008\\[-1pt]{\footnotesize [-0.038, 0.026]}\end{tabular} & \begin{tabular}[c]{@{}c@{}}0.013\\[-1pt]{\footnotesize [-0.018, 0.049]}\end{tabular} & \begin{tabular}[c]{@{}c@{}}-0.057\\[-1pt]{\footnotesize [-0.095, -0.015]}\end{tabular} & \begin{tabular}[c]{@{}c@{}}-0.102\\[-1pt]{\footnotesize [-0.138, -0.063]}\end{tabular} \\
$t_1$의 송금 입차수(기준선) & \begin{tabular}[c]{@{}c@{}}0.198\\[-1pt]{\footnotesize [0.168, 0.228]}\end{tabular} & \begin{tabular}[c]{@{}c@{}}0.169\\[-1pt]{\footnotesize [0.135, 0.199]}\end{tabular} & \begin{tabular}[c]{@{}c@{}}0.016\\[-1pt]{\footnotesize [-0.017, 0.048]}\end{tabular} & \begin{tabular}[c]{@{}c@{}}0.011\\[-1pt]{\footnotesize [-0.017, 0.041]}\end{tabular} \\
\midrule
EndorseRank & \begin{tabular}[c]{@{}c@{}}0.072\\[-1pt]{\footnotesize [0.041, 0.111]}\end{tabular} & \begin{tabular}[c]{@{}c@{}}0.048\\[-1pt]{\footnotesize [0.017, 0.082]}\end{tabular} & \begin{tabular}[c]{@{}c@{}}-0.016\\[-1pt]{\footnotesize [-0.053, 0.024]}\end{tabular} & \begin{tabular}[c]{@{}c@{}}-0.016\\[-1pt]{\footnotesize [-0.054, 0.023]}\end{tabular} \\
AWP & \begin{tabular}[c]{@{}c@{}}0.210\\[-1pt]{\footnotesize [0.181, 0.238]}\end{tabular} & \begin{tabular}[c]{@{}c@{}}0.191\\[-1pt]{\footnotesize [0.159, 0.219]}\end{tabular} & \begin{tabular}[c]{@{}c@{}}0.016\\[-1pt]{\footnotesize [-0.015, 0.044]}\end{tabular} & \begin{tabular}[c]{@{}c@{}}-0.020\\[-1pt]{\footnotesize [-0.051, 0.006]}\end{tabular} \\
\midrule
C-PR ($\lambda=1$) & \begin{tabular}[c]{@{}c@{}}0.073\\[-1pt]{\footnotesize [0.043, 0.111]}\end{tabular} & \begin{tabular}[c]{@{}c@{}}0.048\\[-1pt]{\footnotesize [0.018, 0.082]}\end{tabular} & \begin{tabular}[c]{@{}c@{}}-0.015\\[-1pt]{\footnotesize [-0.051, 0.027]}\end{tabular} & \begin{tabular}[c]{@{}c@{}}-0.016\\[-1pt]{\footnotesize [-0.054, 0.023]}\end{tabular} \\
C-PR ($\lambda=0$) & \begin{tabular}[c]{@{}c@{}}0.131\\[-1pt]{\footnotesize [0.101, 0.162]}\end{tabular} & \begin{tabular}[c]{@{}c@{}}0.153\\[-1pt]{\footnotesize [0.120, 0.182]}\end{tabular} & \begin{tabular}[c]{@{}c@{}}0.029\\[-1pt]{\footnotesize [-0.003, 0.058]}\end{tabular} & \begin{tabular}[c]{@{}c@{}}0.012\\[-1pt]{\footnotesize [-0.014, 0.045]}\end{tabular} \\
C-PR ($\lambda=0.25$) & \begin{tabular}[c]{@{}c@{}}0.128\\[-1pt]{\footnotesize [0.096, 0.158]}\end{tabular} & \begin{tabular}[c]{@{}c@{}}0.154\\[-1pt]{\footnotesize [0.122, 0.184]}\end{tabular} & \begin{tabular}[c]{@{}c@{}}0.025\\[-1pt]{\footnotesize [-0.007, 0.055]}\end{tabular} & \begin{tabular}[c]{@{}c@{}}0.007\\[-1pt]{\footnotesize [-0.019, 0.038]}\end{tabular} \\
C-PR ($\lambda=0.5$) & \begin{tabular}[c]{@{}c@{}}0.121\\[-1pt]{\footnotesize [0.089, 0.152]}\end{tabular} & \begin{tabular}[c]{@{}c@{}}0.151\\[-1pt]{\footnotesize [0.120, 0.183]}\end{tabular} & \begin{tabular}[c]{@{}c@{}}0.024\\[-1pt]{\footnotesize [-0.009, 0.053]}\end{tabular} & \begin{tabular}[c]{@{}c@{}}0.007\\[-1pt]{\footnotesize [-0.019, 0.039]}\end{tabular} \\
C-PR ($\lambda=0.75$) & \begin{tabular}[c]{@{}c@{}}0.111\\[-1pt]{\footnotesize [0.079, 0.140]}\end{tabular} & \begin{tabular}[c]{@{}c@{}}0.146\\[-1pt]{\footnotesize [0.114, 0.177]}\end{tabular} & \begin{tabular}[c]{@{}c@{}}0.023\\[-1pt]{\footnotesize [-0.009, 0.052]}\end{tabular} & \begin{tabular}[c]{@{}c@{}}0.009\\[-1pt]{\footnotesize [-0.019, 0.041]}\end{tabular} \\
S-PR & \begin{tabular}[c]{@{}c@{}}0.141\\[-1pt]{\footnotesize [0.114, 0.173]}\end{tabular} & \begin{tabular}[c]{@{}c@{}}0.120\\[-1pt]{\footnotesize [0.089, 0.150]}\end{tabular} & \begin{tabular}[c]{@{}c@{}}0.014\\[-1pt]{\footnotesize [-0.016, 0.047]}\end{tabular} & \begin{tabular}[c]{@{}c@{}}-0.008\\[-1pt]{\footnotesize [-0.034, 0.021]}\end{tabular} \\
\bottomrule
\end{tabular}
}
\end{table}


송금으로 만든 점수들은 이익을 낸 닫기 횟수와 실현 이익(실제로 손에 넣은 이익)을 줄 세워요. AWP는 $0.210$과 $0.191$, 송금 입차수는 $0.198$과 $0.169$예요. 금액 그대로의 송금 층과 그것을 섞은 점수들을 걷는 점수는 $0.111$에서 $0.141$까지, 그리고 $0.120$에서 $0.154$까지예요. 엔도스랭크와 C-PR ($\lambda=1$)은 약 $0.07$과 $0.05$이고, 받은 허락 수는 아무것도 줄 세우지 못해요. 두 라벨 모두 지갑이 거래를 많이 할수록 커져요. 그렇지 않은 두 비율에서는 $0.03$을 넘는 계수가 없고, 받은 허락 수는 두 비율 모두에서 음수예요. 이것은 활동을 따라가는 송금 점수와, 트레이더 대부분을 동점으로 남기는 허락 점수에서 예상할 만한 결과예요. 어떤 점수도 지갑이 얼마나 거래를 잘하는지는 줄 세우지 못해요. 같은 실행의 Aave 청산 라벨은 지갑 174개만 다루고, 그 가운데 청산된 것은 10개예요. 읽어 내기에는 너무 적어요. 등록한 기간에서는 C-PR에게 따로 트레이더 라벨이 주어져요.

\dissertationSubsubheading[sub:fresh-results]{등록 재현: 2026년 6월부터 8월까지}

\ref{sub:fresh-holdout}절의 등록은 2026년 10월 1일에 공개 저장소에 커밋(기록으로 올림)되었고, 6월부터 8월까지의 로그는 그날 늦게 꺼냈어요. 허락, 송금, GMX 이벤트에 대한 달마다의 질의 아홉 개가 약 1.04\,TB를 훑었고, 허락 로그 163,553개, 송금 로그 1,013,023개, GMX 닫기 144,152개를 돌려주었어요. 이것들은 모두 해독되었어요. 자료 꺼내기 기록표(매니페스트)에는 두 등록 파일의 해시(파일 내용으로 만든 지문 같은 값)가 적혀 있고, 시험용 우회 장치(규칙을 건너뛰게 하는 시험용 스위치)를 쓴 일은 적혀 있지 않아요. 새 숫자를 하나라도 계산하기 전에, 점검으로 평가 코드를 봄 기간에 돌려 보았어요. 이 코드는 표~\ref{tab:holdout-spenders}에도 있는 점수 여덟 개(한 층만 걷는 두 점수, 세 가지 $\lambda$ 값의 C-PR, S-PR, 두 차수)의 봄 계수 열여섯 개와 9월 14일 규칙의 대비를 똑같이 다시 만들어 냈어요. 그리고 위에서 말한, 새 허락 쌍에서 C-PR과 송금 층 사이의 차이 $+0.251$도 내놓았어요. 그다음 등록한 평가를 짝지은 다시 뽑기 400번과 시드 42(무작위 뽑기의 시작 값)로 돌렸어요.

5월 31일 동결일에는 스펜더 1,964개가 0보다 큰 마지막 허락을 가지고 있었어요. 그 가운데 158개가 6월부터 8월까지 새 허락 쌍을 얻었고, 136개가 새 송금자를 얻었어요. 매칭 지갑 5,521개 가운데 5,151개가 동결일 그래프에 있었고, 그 가운데 1,402개가 라벨 기간에 포지션을 적어도 세 번 닫았어요. 1,402개 가운데 789개는 청산을 적어도 한 번 겪었고, 17개는 모든 닫기가 청산이었어요. 표~\ref{tab:fresh-holdout}는 등록한 모든 점수를 등록한 라벨 세 개에서 보여 주고, 표~\ref{tab:fresh-contrasts}는 대비와 그 판정을 보여 줘요. 등록 문서는 한 층만 걷는 두 점수를 엔도스랭크와 AWP라고 불러요. 표에서는 이 논문의 이름을 써요.

% Registered replication, June-August 2026 -- regenerate with:
%   2-Dissertation-Draft-Works/wallet-reputation-experiments/scripts/run_fresh_holdout.py
% Generated by export_latex_results.py -- do not edit by hand unless re-export fails
\begin{table}[htbp]
\centering
\caption[Registered replication: Kendall $\tau$ between scores frozen at 31~May~2026 and the three registered labels, counted from 1~June to 31~August~2026.]{Registered replication: Kendall $\tau$ between scores frozen at 31~May~2026 and the three registered labels, counted from 1~June to 31~August~2026. The spenders are every spender holding a positive latest allowance at the freeze; 158 of them gained a new approval pair and 136 a new sender. The traders are the 5{,}151 matched wallets in the freeze-date graph; the label is defined for the 1{,}402 with at least three closes in the label window, 789 of which had at least one liquidation. C-PR ($\lambda=0$) is the comparator of the registered rule and AWP that of a registered sensitivity analysis (Table~\ref{tab:fresh-contrasts}). CI = 95\% percentile bootstrap over wallets (400 paired resamples, seed 42).}
\label{tab:fresh-holdout}
\fitwidth{%
\begin{tabular}{lccc}
\toprule
 & \multicolumn{2}{c}{Spenders ($n=1{,}964$)} & Traders ($n=1{,}402$) \\
\cmidrule(lr){2-3} \cmidrule(lr){4-4}
Score at 31 May 2026 & New approval pairs & New transfer senders & Liquidation-free close rate \\
\midrule
In-approve degree at the freeze (baseline) & 0.613 [0.564, 0.663] & 0.371 [0.313, 0.429] & 0.144 [0.110, 0.172] \\
Transfer in-degree at the freeze (baseline) & 0.128 [0.064, 0.188] & 0.435 [0.401, 0.468] & -0.006 [-0.046, 0.031] \\
\midrule
AWP & 0.118 [0.071, 0.161] & 0.319 [0.291, 0.345] & 0.045 [0.009, 0.079] \\
\midrule
C-PR ($\lambda=1$) & 0.338 [0.292, 0.388] & 0.289 [0.247, 0.332] & 0.086 [0.037, 0.129] \\
C-PR ($\lambda=0$) & 0.058 [0.012, 0.103] & 0.257 [0.229, 0.286] & 0.065 [0.030, 0.102] \\
C-PR ($\lambda=0.25$) & 0.228 [0.190, 0.265] & 0.245 [0.217, 0.273] & 0.077 [0.042, 0.114] \\
C-PR ($\lambda=0.5$) & 0.238 [0.200, 0.274] & 0.246 [0.218, 0.274] & 0.083 [0.047, 0.119] \\
C-PR ($\lambda=0.75$) & 0.243 [0.206, 0.280] & 0.248 [0.219, 0.275] & 0.086 [0.051, 0.121] \\
S-PR & 0.126 [0.076, 0.172] & 0.254 [0.223, 0.281] & 0.030 [-0.006, 0.066] \\
\bottomrule
\end{tabular}
}
\end{table}


% Registered replication, June-August 2026 -- regenerate with:
%   2-Dissertation-Draft-Works/wallet-reputation-experiments/scripts/run_fresh_holdout.py
% Generated by export_latex_results.py -- do not edit by hand unless re-export fails
\begin{table}[htbp]
\centering
\caption[등록 재현: 짝지은 대비 F1--F6과 그 판정, 그리고 자료를 꺼내기 전에 더한 민감도 분석 두 가지.]{등록 재현의 짝지은 대비 F1--F6과 그 판정, 그리고 자료를 꺼내기 전에 더한 민감도 분석 두 가지예요. $\Delta\tau=\tau_a-\tau_b$이고, C-PR이 $a$예요. 등록 문서에서는 한 층만 걷는 두 점수, 곧 C-PR ($\lambda=1$)과 C-PR ($\lambda=0$)을 엔도스랭크와 AWP라고 불러요. 우월성(더 나음) 시험은 구간이 0보다 위에 있으면 통과해요. 비열등성(허용 폭보다 더 나쁘지 않음) 시험은 구간의 아래 끝이 $-0.02$보다 위에 있으면 통과해요. 9월 14일 규칙의 첫 부분인 F6a는 구간이 0을 품거나 0보다 위에 있으면 통과해요. F1--F3은 표~\ref{tab:fresh-holdout}처럼 스펜더(새 허락 쌍, 새 송금자)와 트레이더(청산 없이 닫은 비율)에서 구했어요. 신뢰구간은 지갑을 다시 뽑아 구한 95\% 백분위 부트스트랩 구간이에요(짝지은 다시 뽑기 400번, 시드 42).}
\label{tab:fresh-contrasts}
\fitwidth{%
\begin{tabular}{lllcccll}
\toprule
 & 대비 & 라벨 & 시험 & $\tau_a$ & $\tau_b$ & $\Delta\tau$ [95\% 신뢰구간] & 판정 \\
\midrule
\multicolumn{8}{l}{\emph{등록한 규칙: C-PR ($\lambda=0.5$) 대 송금 층만 걷는 C-PR ($\lambda=0$)}} \\
F1 & C-PR 대 C-PR ($\lambda=0$) & 새 허락 쌍 & 우월성 & 0.238 & 0.058 & +0.179 [0.132, 0.225] & 통과 \\
F2 & C-PR 대 C-PR ($\lambda=0$) & 새 송금자 & 비열등성, $-0.02$ & 0.246 & 0.257 & -0.011 [-0.025, 0.003] & 실패 \\
F3 & C-PR 대 C-PR ($\lambda=0$) & 청산 없이 닫은 비율 & 비열등성, $-0.02$ & 0.083 & 0.065 & +0.018 [0.011, 0.024] & 통과 \\
\multicolumn{7}{l}{결정: F1, F2, F3이 모두 통과해야 함} & 충족 안 됨 \\
\midrule
\multicolumn{8}{l}{\emph{보고만 하고 결정에는 넣지 않음}} \\
F4 & C-PR 대 C-PR ($\lambda=0$) & 청산 없이 닫은 비율 & 우월성 & 0.083 & 0.065 & +0.018 [0.011, 0.024] & 통과 \\
F5 & C-PR 대 C-PR ($\lambda=1$) & 청산 없이 닫은 비율 & 우월성 & 0.083 & 0.086 & -0.003 [-0.060, 0.051] & 실패 \\
F6a & C-PR 대 C-PR ($\lambda=1$) & 새 허락 쌍 & 더 나쁘지 않음 & 0.238 & 0.338 & -0.101 [-0.127, -0.076] & 실패 \\
F6b & C-PR 대 C-PR ($\lambda=0$) & 새 송금자 & 우월성 & 0.246 & 0.257 & -0.011 [-0.025, 0.003] & 실패 \\
\midrule
\multicolumn{8}{l}{\emph{민감도 분석: AWP를 비교 대상으로}} \\
F1 & C-PR 대 AWP & 새 허락 쌍 & 우월성 & 0.238 & 0.118 & +0.119 [0.075, 0.169] & 통과 \\
F2 & C-PR 대 AWP & 새 송금자 & 비열등성, $-0.02$ & 0.246 & 0.319 & -0.073 [-0.092, -0.055] & 실패 \\
F3 & C-PR 대 AWP & 청산 없이 닫은 비율 & 비열등성, $-0.02$ & 0.083 & 0.045 & +0.037 [0.011, 0.066] & 통과 \\
F4 & C-PR 대 AWP & 청산 없이 닫은 비율 & 우월성 & 0.083 & 0.045 & +0.037 [0.011, 0.066] & 통과 \\
\midrule
\multicolumn{8}{l}{\emph{민감도 분석: 그래프 밖의 지갑을 외톨이 점으로 둠}} \\
F1 & C-PR 대 C-PR ($\lambda=0$) & 새 허락 쌍 & 우월성 & 0.238 & 0.053 & +0.185 [0.137, 0.232] & 통과 \\
F2 & C-PR 대 C-PR ($\lambda=0$) & 새 송금자 & 비열등성, $-0.02$ & 0.246 & 0.258 & -0.012 [-0.027, 0.003] & 실패 \\
F3 & C-PR 대 C-PR ($\lambda=0$) & 청산 없이 닫은 비율 & 비열등성, $-0.02$ & 0.083 & 0.065 & +0.018 [0.010, 0.024] & 통과 \\
F4 & C-PR 대 C-PR ($\lambda=0$) & 청산 없이 닫은 비율 & 우월성 & 0.083 & 0.065 & +0.018 [0.010, 0.024] & 통과 \\
F5 & C-PR 대 C-PR ($\lambda=1$) & 청산 없이 닫은 비율 & 우월성 & 0.083 & 0.155 & -0.072 [-0.118, -0.028] & 실패 \\
F6a & C-PR 대 C-PR ($\lambda=1$) & 새 허락 쌍 & 더 나쁘지 않음 & 0.238 & 0.338 & -0.101 [-0.127, -0.076] & 실패 \\
F6b & C-PR 대 C-PR ($\lambda=0$) & 새 송금자 & 우월성 & 0.246 & 0.258 & -0.012 [-0.027, 0.003] & 실패 \\
\bottomrule
\end{tabular}
}
\end{table}


규칙은 충족되지 않아요. F1은 성립해요. C-PR은 나중에 새 허락 쌍을 얻은 스펜더들을 $\tau=0.238$로 줄 세우고, 송금 층만 걷는 점수는 $0.058$이에요. 차이는 $+0.179$예요(구간 $[0.132,0.225]$). 그래서 권한 주기(허락)에서 봄에 보인 이득이 다시 나타나요. F3도 성립해요. 트레이더의 청산 없이 닫은 비율(닫은 포지션 가운데 청산되지 않은 것의 비율)에서 C-PR은 $0.083$이고, 송금 걷기는 $0.065$예요($+0.018$, $[0.011,0.024]$). F2는 성립하지 않아요. 새 송금자에서 C-PR은 $0.246$이고 송금 걷기는 $0.257$이에요. 차이는 $-0.011$이고, 구간의 아래 끝 $-0.025$가 허용 폭(이만큼까지는 뒤져도 괜찮다고 미리 정한 폭) $-0.02$보다 낮아요. 구간에는 0도 들어 있어요. 그래서 이 자료는 새 송금자에서 C-PR이 송금 걷기보다 못하다는 것을 보여 주지 않아요. 그렇다고 둘의 차이가 $0.02$ 안쪽이라는 것을 보여 주지도 않는데, 규칙이 요구한 것은 바로 그것이었어요. 규칙에는 세 가설이 모두 필요했어요. 그래서 계획에서 이런 경우에 대해 말한 대로, 이 논문은 C-PR이 송금만 걷는 점수보다 낫다고 주장하지 않아요.

계획은 두 계수가 모두 0에 가까워도 F3이 통과할 수 있다고 미리 경고했어요. 두 계수는 모두 작지만, 두 구간 모두 0보다 위에 있어요(C-PR $[0.047,0.119]$, 송금 걷기 $[0.030,0.102]$). 그래서 두 걷기는 저마다 어떤 트레이더가 청산을 피하는지에 대해 정보를 조금 담고 있어요. 동결일의 받은 허락 수는 이 라벨을 어떤 걷기보다도 잘 줄 세워요($0.144$, $[0.110,0.172]$). 허락 걷기는 $0.086$, AWP는 $0.045$예요.

결정 옆에 함께 보고하는 결과들도 같은 순서를 따라요. F4는 성립해요. 청산 없이 닫은 비율에서 C-PR은 송금 걷기보다 못하지 않은 데 그치지 않고, 더 나아요. F5는 성립하지 않아요. 그 라벨에서 C-PR과 허락 걷기를 구별할 수 없기 때문이에요($-0.003$, $[-0.060,0.051]$). 9월 14일 규칙의 주된 부분인 F6은 이번에도 두 부분 모두 통과하지 못해요. C-PR은 새 허락 쌍에서 허락 걷기보다 낮고($-0.101$, $[-0.127,-0.076]$), 새 송금자에서 송금 걷기보다 높지 않아요. 봄처럼, C-PR의 계수는 두 스펜더 라벨 모두에서 $\lambda$와 함께 올라가요. S-PR은 새 송금자에서 송금 걷기와 같은 정도이고($-0.003$, $[-0.017,0.010]$), 새 허락 쌍에서는 허락 걷기보다 한참 낮아요($-0.213$).

등록한 민감도 분석 두 가지는 어느 것도 판정을 바꾸지 않아요(표~\ref{tab:fresh-contrasts}). AWP를 비교 대상으로 쓰면, F1과 F3은 여전히 성립하고($+0.119$와 $+0.037$), F2는 뚜렷하게 성립하지 않아요. AWP는 새 송금자를 $0.319$로 줄 세우고, AWP와 견준 C-PR의 차이는 $-0.073$이에요($[-0.092,-0.055]$). 이 구간은 통째로 허용 폭보다 아래에 있어요. 그러니까 AWP와 견주면, C-PR은 나중 허락에서 얻는 이득과 맞바꾸어, 나중에 들어오는 흐름을 줄 세우는 능력의 잴 수 있을 만큼을 내줘요. 그래프 밖의 지갑을 외톨이 점으로 남겨 두면, F1부터 F4까지는 판정이 그대로예요. F5는 ‘차이 없음’에서 ‘손해’로 바뀌어요. 이 규칙은 허락 그래프 밖에 있는 매칭 트레이더들을, 허락 걷기 순서의 바닥에서 들어오는 화살표가 없는 지갑들의 덩어리로 옮겨요(\ref{sec:rank-divergence}절). 그러면 청산 없이 닫은 비율에서 허락 걷기의 계수가 $0.086$에서 $0.155$로 올라가고, C-PR은 그보다 낮아져요($-0.072$, $[-0.118,-0.028]$). F6은 이 규칙에서도 통과하지 못해요.

% Registered replication, June-August 2026 -- regenerate with:
%   2-Dissertation-Draft-Works/wallet-reputation-experiments/scripts/run_fresh_posthoc.py
% Generated by export_latex_results.py -- do not edit by hand unless re-export fails
\begin{table}[htbp]
\centering
\caption[EndorseRank on the cohorts and labels of the registered replication, scored after that replication was evaluated; none of these rows enters its decision.]{EndorseRank on the cohorts and labels of the registered replication, scored after that replication was evaluated; none of these rows enters its decision. The AWP, C-PR and degree rows repeat Table~\ref{tab:fresh-holdout}, and the contrasts use the same wallet resamples; the differences between EndorseRank and AWP are in Table~\ref{tab:endorserank-awp-holdout}. CI = 95\% percentile bootstrap over wallets (400 paired resamples, seed 42).}
\label{tab:fresh-posthoc}
\fitwidth{%
\begin{tabular}{lccc}
\toprule
 & \multicolumn{2}{c}{Spenders ($n=1{,}964$)} & Traders \\
\cmidrule(lr){2-3} \cmidrule(lr){4-4}
Score at 31 May 2026 & New approval pairs & New transfer senders & Liquidation-free close rate \\
\midrule
In-approve degree at the freeze (baseline) & 0.613 [0.564, 0.663] & 0.371 [0.313, 0.429] & 0.144 [0.110, 0.172] \\
Transfer in-degree at the freeze (baseline) & 0.128 [0.064, 0.188] & 0.435 [0.401, 0.468] & -0.006 [-0.046, 0.031] \\
\midrule
EndorseRank & 0.333 [0.306, 0.364] & 0.255 [0.227, 0.284] & 0.081 [0.032, 0.123] \\
EndorseRank, restarts by approving activity & 0.325 [0.295, 0.358] & 0.284 [0.252, 0.312] & 0.077 [0.037, 0.116] \\
AWP & 0.118 [0.071, 0.161] & 0.319 [0.291, 0.345] & 0.045 [0.009, 0.079] \\
C-PR ($\lambda=0.5$) & 0.238 [0.200, 0.274] & 0.246 [0.218, 0.274] & 0.083 [0.047, 0.119] \\
\midrule
EndorseRank minus t1 in-approve degree, new approval pairs & \multicolumn{3}{l}{$\Delta\tau=-0.280$ [-0.313, -0.247]} \\
AWP minus t1 in-degree, new transfer senders & \multicolumn{3}{l}{$\Delta\tau=-0.116$ [-0.133, -0.095]} \\
C-PR minus EndorseRank, new approval pairs & \multicolumn{3}{l}{$\Delta\tau=-0.096$ [-0.128, -0.065]} \\
C-PR minus EndorseRank, liquidation-free close rate & \multicolumn{3}{l}{$\Delta\tau=+0.002$ [-0.053, 0.056]} \\
\bottomrule
\end{tabular}
}
\end{table}


엔도스랭크는 등록한 점수에 들어 있지 않았어요. 그래서 표~\ref{tab:fresh-posthoc}은 등록한 평가가 끝난 뒤에, 같은 코호트, 같은 라벨, 같은 지갑 다시 뽑기로 엔도스랭크의 점수를 매겨요. 그 줄들은 어느 것도 결정에 들어가지 않아요. 봄의 모습은 이 기간에서도 모든 페이지랭크에서 그대로 나타나요. 동결일의 받은 허락 수는 이번에도 새 허락 쌍을 어떤 걷기보다도 잘 줄 세워요($0.613$; 엔도스랭크에서 차수를 뺀 값 $-0.280$, $[-0.313,-0.247]$). 송금 입차수도 이번에도 새 송금자를 어떤 걷기보다도 잘 줄 세워요($0.435$; AWP에서 차수를 뺀 값 $-0.116$, $[-0.133,-0.095]$). C-PR은 이번에도 새 허락 쌍에서 엔도스랭크보다 낮아요($-0.096$, $[-0.128,-0.065]$). 트레이더의 청산 없이 닫은 비율에서 엔도스랭크는 $0.081$로, C-PR과 비슷해요($0.083$; $+0.002$, $[-0.053,0.056]$). AWP의 재시작을 쓰면, 엔도스랭크는 두 스펜더 라벨에서 $0.325$와 $0.284$, 트레이더에서 $0.077$이 나와요.

% Registered replication, June-August 2026 -- regenerate with:
%   2-Dissertation-Draft-Works/wallet-reputation-experiments/scripts/run_fresh_holdout.py
% Generated by export_latex_results.py -- do not edit by hand unless re-export fails
\begin{table}[htbp]
\centering
\caption[Registered replication, traders: Kendall $\tau$ between scores frozen at 31~May~2026 and the other GMX~V2 labels counted from June to August~2026, on the 1{,}402 matched wallets with at least three closes in that window.]{Registered replication, traders: Kendall $\tau$ between scores frozen at 31~May~2026 and the other GMX~V2 labels counted from June to August~2026, on the 1{,}402 matched wallets with at least three closes in that window. No liquidation is 1 when none of the wallet's closes was a liquidation. Profitable closes and realized gain grow with the number of closes; the two shares do not. These labels are reported and do not enter the decision. CI = 95\% percentile bootstrap over wallets (400 paired resamples, seed 42).}
\label{tab:fresh-traders}
\fitwidth{%
\begin{tabular}{lccccc}
\toprule
Score at 31 May 2026 & No liquidation & Profitable closes & Realized gain & Profitable share & Non-loss share \\
\midrule
In-approve degree at the freeze (baseline) & \begin{tabular}[c]{@{}c@{}}0.187\\[-1pt]{\footnotesize [0.150, 0.221]}\end{tabular} & \begin{tabular}[c]{@{}c@{}}-0.004\\[-1pt]{\footnotesize [-0.033, 0.028]}\end{tabular} & \begin{tabular}[c]{@{}c@{}}0.021\\[-1pt]{\footnotesize [-0.015, 0.054]}\end{tabular} & \begin{tabular}[c]{@{}c@{}}-0.074\\[-1pt]{\footnotesize [-0.105, -0.042]}\end{tabular} & \begin{tabular}[c]{@{}c@{}}-0.139\\[-1pt]{\footnotesize [-0.172, -0.101]}\end{tabular} \\
Transfer in-degree at the freeze (baseline) & \begin{tabular}[c]{@{}c@{}}-0.072\\[-1pt]{\footnotesize [-0.115, -0.030]}\end{tabular} & \begin{tabular}[c]{@{}c@{}}0.191\\[-1pt]{\footnotesize [0.150, 0.224]}\end{tabular} & \begin{tabular}[c]{@{}c@{}}0.127\\[-1pt]{\footnotesize [0.089, 0.166]}\end{tabular} & \begin{tabular}[c]{@{}c@{}}0.026\\[-1pt]{\footnotesize [-0.008, 0.060]}\end{tabular} & \begin{tabular}[c]{@{}c@{}}-0.001\\[-1pt]{\footnotesize [-0.036, 0.038]}\end{tabular} \\
\midrule
AWP & \begin{tabular}[c]{@{}c@{}}-0.030\\[-1pt]{\footnotesize [-0.069, 0.011]}\end{tabular} & \begin{tabular}[c]{@{}c@{}}0.233\\[-1pt]{\footnotesize [0.193, 0.268]}\end{tabular} & \begin{tabular}[c]{@{}c@{}}0.182\\[-1pt]{\footnotesize [0.142, 0.217]}\end{tabular} & \begin{tabular}[c]{@{}c@{}}0.030\\[-1pt]{\footnotesize [-0.004, 0.060]}\end{tabular} & \begin{tabular}[c]{@{}c@{}}-0.035\\[-1pt]{\footnotesize [-0.065, -0.001]}\end{tabular} \\
\midrule
C-PR ($\lambda=1$) & \begin{tabular}[c]{@{}c@{}}0.069\\[-1pt]{\footnotesize [0.011, 0.116]}\end{tabular} & \begin{tabular}[c]{@{}c@{}}0.074\\[-1pt]{\footnotesize [0.036, 0.107]}\end{tabular} & \begin{tabular}[c]{@{}c@{}}0.016\\[-1pt]{\footnotesize [-0.024, 0.059]}\end{tabular} & \begin{tabular}[c]{@{}c@{}}-0.043\\[-1pt]{\footnotesize [-0.084, -0.000]}\end{tabular} & \begin{tabular}[c]{@{}c@{}}-0.130\\[-1pt]{\footnotesize [-0.166, -0.091]}\end{tabular} \\
C-PR ($\lambda=0$) & \begin{tabular}[c]{@{}c@{}}0.014\\[-1pt]{\footnotesize [-0.026, 0.056]}\end{tabular} & \begin{tabular}[c]{@{}c@{}}0.140\\[-1pt]{\footnotesize [0.103, 0.177]}\end{tabular} & \begin{tabular}[c]{@{}c@{}}0.145\\[-1pt]{\footnotesize [0.110, 0.180]}\end{tabular} & \begin{tabular}[c]{@{}c@{}}0.033\\[-1pt]{\footnotesize [-0.001, 0.069]}\end{tabular} & \begin{tabular}[c]{@{}c@{}}-0.004\\[-1pt]{\footnotesize [-0.040, 0.035]}\end{tabular} \\
C-PR ($\lambda=0.25$) & \begin{tabular}[c]{@{}c@{}}0.032\\[-1pt]{\footnotesize [-0.007, 0.072]}\end{tabular} & \begin{tabular}[c]{@{}c@{}}0.133\\[-1pt]{\footnotesize [0.094, 0.168]}\end{tabular} & \begin{tabular}[c]{@{}c@{}}0.145\\[-1pt]{\footnotesize [0.110, 0.182]}\end{tabular} & \begin{tabular}[c]{@{}c@{}}0.027\\[-1pt]{\footnotesize [-0.007, 0.062]}\end{tabular} & \begin{tabular}[c]{@{}c@{}}-0.011\\[-1pt]{\footnotesize [-0.047, 0.028]}\end{tabular} \\
C-PR ($\lambda=0.5$) & \begin{tabular}[c]{@{}c@{}}0.043\\[-1pt]{\footnotesize [0.004, 0.082]}\end{tabular} & \begin{tabular}[c]{@{}c@{}}0.124\\[-1pt]{\footnotesize [0.084, 0.160]}\end{tabular} & \begin{tabular}[c]{@{}c@{}}0.142\\[-1pt]{\footnotesize [0.107, 0.179]}\end{tabular} & \begin{tabular}[c]{@{}c@{}}0.024\\[-1pt]{\footnotesize [-0.010, 0.058]}\end{tabular} & \begin{tabular}[c]{@{}c@{}}-0.010\\[-1pt]{\footnotesize [-0.045, 0.028]}\end{tabular} \\
C-PR ($\lambda=0.75$) & \begin{tabular}[c]{@{}c@{}}0.053\\[-1pt]{\footnotesize [0.011, 0.091]}\end{tabular} & \begin{tabular}[c]{@{}c@{}}0.110\\[-1pt]{\footnotesize [0.068, 0.147]}\end{tabular} & \begin{tabular}[c]{@{}c@{}}0.135\\[-1pt]{\footnotesize [0.101, 0.171]}\end{tabular} & \begin{tabular}[c]{@{}c@{}}0.021\\[-1pt]{\footnotesize [-0.014, 0.056]}\end{tabular} & \begin{tabular}[c]{@{}c@{}}-0.008\\[-1pt]{\footnotesize [-0.041, 0.030]}\end{tabular} \\
S-PR & \begin{tabular}[c]{@{}c@{}}-0.030\\[-1pt]{\footnotesize [-0.070, 0.014]}\end{tabular} & \begin{tabular}[c]{@{}c@{}}0.176\\[-1pt]{\footnotesize [0.139, 0.211]}\end{tabular} & \begin{tabular}[c]{@{}c@{}}0.132\\[-1pt]{\footnotesize [0.094, 0.170]}\end{tabular} & \begin{tabular}[c]{@{}c@{}}0.040\\[-1pt]{\footnotesize [0.003, 0.075]}\end{tabular} & \begin{tabular}[c]{@{}c@{}}-0.023\\[-1pt]{\footnotesize [-0.057, 0.015]}\end{tabular} \\
\bottomrule
\end{tabular}
}
\end{table}


표~\ref{tab:fresh-traders}은 결정에 들어가지 않는 등록 트레이더 라벨들을 보여 주는데, 이 결과는 봄의 트레이더 실행을 되풀이해요. 등록한 점수 가운데 이익을 낸 닫기 횟수($0.233$)와 실현 이익($0.182$)을 가장 높게 줄 세우는 것은 AWP예요. 송금 입차수는 $0.191$과 $0.127$, C-PR은 $0.124$와 $0.142$예요. 모든 점수가 이익을 낸 닫기의 비율에서는 $0.04$ 이하에, 손실이 아니었던 닫기의 비율에서는 0 이하에 머물러요.

\dissertationSubheading[sec:empirical-claims]{주요 주장(C1--C5)}

이 장의 결과는 아래의 주요 주장들을 뒷받침해요. 각 주장은 표~\ref{tab:hrq-claim-map}의 가설 열, 연구 질문 열과 서로 찾아볼 수 있게 이어져 있어요.

\begin{enumerate}
\item[C1.] 똑같은 지갑 묶음에서, 각 점수의 풀이 프로그램 호출 비용은 그 그래프의 화살표 수를 따라가요. 행렬 조립과 반복이 둘 다 $|E|$에 따라 커지기 때문이에요. 엔도스랭크의 그래프는 화살표가 송금 그래프의 약 9분의 1이고, 호출에 드는 시간과 메모리도 AWP의 그만큼이에요(표~\ref{tab:benchmark-runtime}). 이 비율은 코호트 안 단계와 풀 단계에서 모두 유지돼요(표~\ref{tab:benchmark-scaling}와~\ref{tab:benchmark-scaling-incohort}). 다만 풀 단계는 코호트 그래프보다 화살표가 많은 경우를 시험하지 않아요. 싼 비용은 엔도스랭크만의 것이에요. C-PR과 S-PR은 송금 층을 걷고, AWP보다 오래 걸려요.
\item[C2.] 같은 기간에서 각 점수는 자기가 걷는 화살표 종류의 대리 지표와 맞고, 그 구간에는 0이 들어 있지 않아요. 섞은 점수들은 더 빽빽한 송금 층을 따라가요(표~\ref{tab:alignment-hybrid}). 이 시험들은 안에서의 일관성을 재고, 일치의 일부는 처음부터 생기게 되어 있어요. 외적 타당도(바깥 자료로도 맞는지)는 이 시험들이 다루는 범위 밖이에요. 엔도스랭크의 계수는 허락 그래프 밖의 지갑을 다루는 규칙이 정해요(표~\ref{tab:tie-sensitivity}).
\item[C3.] 엔도스랭크와 AWP는 코호트를 다르게 줄 세워요. 두 점수의 일치는 구간에 0도 1도 들어 있지 않아요. 이 일치는 주로 허락 그래프 밖의 지갑들에서 생기고, 그 지갑들을 외톨이 점으로 남겨 두면 0보다 조금 아래로 떨어져요(표~\ref{tab:tie-sensitivity}, 그림~\ref{fig:rank-scatter}).
\item[C4.] 기간 밖에서는 두 기간 어디에서도, 이 논문의 어떤 페이지랭크도 스펜더 라벨을 그 라벨과 같은 화살표 종류의 원시 차수만큼 잘 줄 세우지 못해요. 차수와의 짝지은 차이는 모두 0보다 아래에 있어요(표~\ref{tab:holdout-diff}와~\ref{tab:fresh-posthoc}). 그리고 이렇게 못 미치는 것은 포함 범위 때문이 아니에요(표~\ref{tab:holdout-receiving}). 매칭 트레이더들에서는 송금으로 만든 점수들이 나중에 이익을 낸 닫기 횟수를 줄 세우지만, 이익을 낸 닫기의 비율을 줄 세우는 점수는 없어요(표~\ref{tab:holdout-traders}). C-PR ($\lambda=1$)의 증가분(원시 차수보다 나은 정도)을 빼면, 이 비교들은 어느 것도 미리 정한 것이 아니에요.
\item[C5.] 두 화살표 종류를 함께 걷는 무작위 걷기 하나는, 한 층만 걷는 두 점수를 이기지 못해요. 그리고 등록 재현은 그것이 송금 걷기보다 낫다는 것을 확인해 주지 않아요. 9월 14일의 규칙에서, C-PR은 시험한 모든 $\lambda$에서 새 허락 쌍에 대해 두 걷기 사이에 놓이고, S-PR은 송금 걷기와 같은 결과를 내요(표~\ref{tab:holdout-spenders}와~\ref{tab:holdout-diff}). 라벨을 본 뒤에 찾은, 새 허락 쌍에서 송금 걷기보다 나았던 봄의 이득은 등록한 기간에서 다시 나타나요(F1, 그리고 F3과 F4). 하지만 새 송금자에서의 차이는 $0.02$ 안쪽으로 묶을 수 없어요(F2). 그래서 등록한 규칙은 충족되지 않아요. 등록한 민감도 분석에서 AWP와 견주면, C-PR은 새 송금자에서 뚜렷하게 져요(표~\ref{tab:fresh-contrasts}). 혼합 점수들이 같은 기간 안에서 어디쯤 있는지는 \ref{sec:coupled-results}절에서 보고해요.
\end{enumerate}

순위 갈라짐은 RQ1에, 같은 기간 표들은 RQ2에, 벤치마크 표들은 RQ3에, 홀드아웃들은 RQ4와 RQ5에 답해요. 엔도스랭크와 AWP를 직접 견준 결과는 부록~\ref{ch:appendix-er-awp}에 있고, 그 가운데 어느 것도 이 논문의 주장을 담고 있지 않아요. \ref{ch:discussion}장은 이 결과들을 \ref{ch:introduction}장의 연구 질문과 목표에 비추어 읽어요.
